\documentclass[11pt]{article}
\input{../shared/project_style.tex}
\rhead{Module 1}
\begin{document}

\DocTitle{Module 1: Ballistic-Diffusive Heat Conduction}{From microscopic carrier transport to a FEM-ready nanoscale heat-transport equation}{Physics note for FEM\_BallisticDiffusion\_PINN}

\begin{overviewbox}
\textbf{Purpose.} This note develops the ballistic-diffusive heat-conduction model from first principles. The goal is not only to write down the final partial differential equation, but to understand why it has the form that it does. We begin with conservation of energy and the microscopic carrier picture, derive the Boltzmann transport equation under the relaxation-time approximation, recover Fourier conduction as a limiting approximation, explain why Fourier and Cattaneo models fail in the quasi-ballistic regime, and then derive the ballistic-diffusive equation used in this project. The final sections translate the physics into a weak form appropriate for finite element method (FEM) implementation.
\end{overviewbox}

\ProjectTOC

\SymbolsAcronymsSection
\begin{symtable}
BTE & Boltzmann transport equation & Phase-space transport equation for heat carriers such as phonons, electrons, or molecules. \\
BD & Ballistic-diffusive & Reduced model in which boundary-originating carriers are treated ballistically and internally scattered carriers are treated diffusively. \\
FEM & Finite element method & Numerical method based on a weak form, basis functions, and element-wise integration. \\
PDE & Partial differential equation & Equation involving partial derivatives with respect to space and/or time. \\
RTA & Relaxation-time approximation & Collision model in which the distribution relaxes toward local equilibrium over a characteristic time. \\
P1 approximation & First-order spherical harmonics approximation & Angular-moment closure that keeps the isotropic part and first directional anisotropy of a distribution. \\
\(t\) & Time & Independent time coordinate. \\
\(\rvec\) & Position vector & Spatial coordinate in the material domain. \\
\(x\) & One-dimensional coordinate & Spatial coordinate across a thin film. \\
\(L\) & Characteristic length & Film thickness or representative device length scale. \\
\(s\) & Ray coordinate & Distance measured along a carrier propagation direction. \\
\(s_0\) & Boundary ray coordinate & Ray coordinate at the boundary point from which a carrier originates. \\
\(\hat{\Omega}\) & Propagation direction & Unit vector in the direction of carrier motion. \\
\(\bm{v}\) & Carrier group velocity & Velocity vector of a carrier wave packet. \\
\(v\) & Carrier speed & Magnitude of the carrier group velocity, \(v=\abs{\bm{v}}\). \\
\(\mu\) & Direction cosine & In one-dimensional transport, \(\mu=\cos\theta\), where \(\theta\) is the angle relative to the positive \(x\)-axis. \\
\(\omega\) & Angular frequency & Carrier angular frequency. \\
\(\hbar\) & Reduced Planck constant & Quantum of action divided by \(2\pi\). \\
\(D(\omega)\) & Density of states & Number of available carrier modes per unit volume per angular-frequency interval. \\
\(f\) & Distribution function & Carrier occupation in phase space. \\
\(f_0\) & Local-equilibrium distribution & Bose-Einstein or Fermi-Dirac distribution determined by local thermodynamic variables. \\
\(f_b\) & Ballistic distribution & Boundary-originating portion of \(f\). \\
\(f_m\) & Medium/scattered distribution & Portion of \(f\) associated with carriers scattered or emitted inside the medium. \\
\(I_\omega\) & Spectral intensity & Directional energy-flow representation of the distribution function. \\
\(I_{0\omega}\) & Equilibrium spectral intensity & Intensity corresponding to local equilibrium. \\
\(I_{b\omega}\) & Ballistic intensity & Boundary-originating part of \(I_\omega\). \\
\(I_{m\omega}\) & Medium/scattered intensity & Internally scattered or emitted part of \(I_\omega\). \\
\(I_{w\omega}\) & Wall/boundary intensity & Intensity emitted or reflected from a boundary into the domain. \\
\(J_{0\omega}\) & Isotropic intensity coefficient & Zeroth angular harmonic of \(I_{m\omega}\). \\
\(\bm{J}_{1\omega}\) & First angular coefficient & First angular harmonic vector of \(I_{m\omega}\). \\
\(\tau_\omega\) & Mode-dependent relaxation time & Average time between scattering events for carriers of frequency \(\omega\). \\
\(\tau\) & Averaged relaxation time & Spectrally averaged relaxation time used in the reduced model. \\
\(\mfp_\omega\) & Mode-dependent mean free path & \(\mfp_\omega=v\tau_\omega\). \\
\(\mfp\) & Mean free path & Average distance traveled between scattering events. \\
\(\mathrm{Kn}\) & Knudsen number & \(\mathrm{Kn}=\mfp/L\), ratio of mean free path to device length. \\
\(u\) & Internal energy density & Carrier energy per unit volume. \\
\(u_b\) & Ballistic internal energy density & Internal energy density associated with \(I_{b\omega}\). \\
\(u_m\) & Medium/scattered internal energy density & Internal energy density associated with \(I_{m\omega}\). \\
\(u_{b0}\) & Equilibrium ballistic energy density & Ballistic energy contribution in a reference equilibrium state. \\
\(u_{m0}\) & Equilibrium medium energy density & Medium energy contribution in a reference equilibrium state. \\
\(T\) & Equivalent temperature & Temperature inferred from internal energy; a true thermodynamic temperature only when local equilibrium is valid. \\
\(T_m\) & Medium-equivalent temperature & \(T_m=u_m/C\) when constant volumetric heat capacity is assumed. \\
\(T_b\) & Ballistic-equivalent temperature & \(T_b=u_b/C\) when constant volumetric heat capacity is assumed. \\
\(T_0\) & Initial/reference temperature & Uniform initial temperature of a validation film. \\
\(T_1\) & Emitted boundary temperature & Temperature characterizing carriers emitted from the hot boundary. \\
\(\Delta T\) & Temperature scale & \(\Delta T=T_1-T_0\). \\
\(C\) & Volumetric heat capacity & \(C=\partial u/\partial T\), energy per unit volume per kelvin. \\
\(C_\omega\) & Spectral heat capacity & Heat capacity contribution from carriers near angular frequency \(\omega\). \\
\(k\) & Thermal conductivity & Diffusive thermal conductivity. \\
\(\alpha\) & Thermal diffusivity & \(\alpha=k/C\). \\
\(\qvec\) & Total heat flux & Energy flow per unit area per unit time. \\
\(\qvec_b\) & Ballistic heat flux & Heat flux carried by boundary-originating carriers. \\
\(\qvec_m\) & Medium/scattered heat flux & Heat flux carried by internally scattered carriers. \\
\(\dot{q}_e\) & Volumetric heat generation & External heat source per unit volume. \\
\(\nvec\) & Unit normal & Outward unit normal vector on a boundary. \\
\(\Omega\) & Spatial domain & Physical region occupied by the simulated material. \\
\(\partial\Omega\) & Boundary of the domain & Surface enclosing \(\Omega\). \\
\(w\) & Test function & Weighting function used to form the FEM weak statement. \\
\(N_i\) & FEM basis function & Shape function associated with node \(i\). \\
\(\phi\) & Generic FEM scalar unknown & In the FEM derivation, \(\phi\) denotes either \(u_m\) or the nondimensional medium energy \(\theta_m\). \\
\(\phi_h\) & FEM approximation & Piecewise-polynomial approximation to \(\phi\). \\
\(\Phi_i\) & Nodal degree of freedom & Coefficient multiplying basis function \(N_i\); usually the value of \(\phi_h\) at node \(i\). \\
\(\Omega_e\) & Element domain & Subregion of the mesh used for element-level integration. \\
\(h_e\) & Element size & Length of a one-dimensional element or characteristic size of a multidimensional element. \\
\(\Gamma_D\) & Dirichlet boundary & Boundary on which \(\phi\) is prescribed. \\
\(\Gamma_N\) & Neumann boundary & Boundary on which normal diffusive flux is prescribed. \\
\(\Gamma_R\) & Robin/Marshak boundary & Boundary with mixed value-gradient relation such as the ballistic-diffusive medium boundary condition. \\
\(\bm{N}\) & Shape-function row vector & Vector collecting element shape functions. \\
\(\bm{B}\) & Gradient matrix & Matrix collecting spatial gradients of shape functions. \\
\(\bm{J}\) & Geometric Jacobian & Matrix mapping natural element coordinates to physical coordinates. \\
\(\xi\) & Natural coordinate & Reference-element coordinate used for element integration. \\
\(M_e,K_e,F_e\) & Element matrices/vectors & Element mass matrix, stiffness matrix, and load/source vector. \\
\(A_{\mathrm{eff}}\) & Effective solve matrix & Matrix solved at one time step after time discretization. \\
\(\Delta t\) & Time step & Increment between two discrete time levels. \\
\(M\) & Mass matrix & FEM matrix from \(\int_\Omega N_i N_j\,\dd V\). \\
\(K\) & Stiffness matrix & FEM matrix from \(\int_\Omega \nabla N_i\cdot \alpha \nabla N_j\,\dd V\). \\
\(\eta\) & Nondimensional coordinate & \(\eta=x/L\). \\
\(t^*\) & Nondimensional time & \(t^*=t/\tau\). \\
\(\theta_m\) & Nondimensional medium energy & Normalized version of \(u_m-u_{m0}\). \\
\(\theta_b\) & Nondimensional ballistic energy & Normalized version of \(u_b-u_{b0}\). \\
\(\theta\) & Nondimensional total energy temperature & Normalized total energy perturbation, \(\theta=\theta_m+\theta_b\). \\
\(q_{b0}\) & Equilibrium ballistic heat flux & Reference ballistic flux subtracted in the normalized validation problem. \\
\(q_b^*\) & Nondimensional ballistic flux & Normalized version of \(q_b-q_{b0}\). \\
\(E_n(z)\) & Exponential integral & Integral \(E_n(z)=\int_1^\infty e^{-zt}t^{-n}\,\dd t\), useful for equilibrium ballistic terms. \\
\end{symtable}

\section{What problem this module is solving}
A superconducting qubit chip contains structures whose relevant dimensions can range from millimeters, to micrometers, to nanometers. A local energy deposition event, a photon absorption event, or a transient heat pulse does not always relax according to ordinary macroscopic Fourier conduction. At small length scales, heat carriers can travel a significant fraction of the device size before scattering. In that case, the heat flux at a point is partly determined by distant boundaries and by the history of carrier flight, not only by the local temperature gradient.

The purpose of this module is to construct a thermal transport model that is more physical than Fourier conduction but much cheaper than a full Boltzmann transport solve. The model should be capable of representing the transition between these limits:
\begin{align}
\text{diffusive limit:} \quad & \mfp \ll L, \\
\text{quasi-ballistic limit:} \quad & \mfp \sim L, \\
\text{ballistic limit:} \quad & \mfp \gtrsim L,
\end{align}
where \(\mfp\) is the carrier mean free path and \(L\) is the characteristic device length.

\begin{physicsbox}
\textbf{Core intuition.} Fourier's law is a local law: the heat flux at \(\rvec\) is determined by \(\nabla T(\rvec)\). Ballistic transport is nonlocal: a carrier at \(\rvec\) may have originated at a boundary point many mean-free-path-scale distances away and may remember the boundary state at an earlier time.
\end{physicsbox}

\section{Heat carriers and the meaning of temperature}
\subsection{Heat is carried by microscopic excitations}
In a solid, thermal energy is transported by microscopic excitations. In an insulating crystal, phonons are usually the dominant heat carriers. In a metal, both electrons and phonons can contribute. In a superconducting qubit material stack, the situation becomes more complicated at cryogenic temperature because phonons, quasiparticles, Cooper pairs, normal-metal traps, and interfaces can all matter. Module 1 intentionally starts with the room-temperature or normal-carrier thermal transport problem because it provides the transport language needed later.

A carrier is described by its position \(\rvec\), propagation direction \(\hat{\Omega}\), frequency \(\omega\), and time \(t\). A distribution function tells us how many carriers occupy a small region of this phase space.

\subsection{Temperature is simple only near local equilibrium}
Macroscopic heat conduction often begins by assuming that every small volume element is in local thermal equilibrium. Then one can assign a temperature \(T(\rvec,t)\), and the local internal energy density is a function of temperature:
\begin{equation}
    u(\rvec,t)=u[T(\rvec,t)].
\end{equation}
For small temperature changes around a reference temperature, this becomes
\begin{equation}
    \dd u = C\,\dd T,
\end{equation}
where \(C\) is volumetric heat capacity. This relation is useful but should not be mistaken for a definition that is always valid. If the carrier distribution is strongly nonequilibrium, a single thermodynamic temperature may not exist.

In this module we use \emph{equivalent temperature} in a precise sense: it is the temperature that an equilibrium distribution would need to have in order to contain the same local internal energy. Therefore, in ballistic transport, temperature is best interpreted as an energy label, not always as a true local-equilibrium state variable.

\section{Macroscopic energy conservation}
Before discussing any constitutive law, we start with energy conservation. Consider an arbitrary fixed control volume \(V\) with boundary \(\partial V\) and outward normal \(\nvec\). Let \(u\) be internal energy density and \(\qvec\) be heat flux. Conservation of energy says
\begin{equation}
    \frac{\dd}{\dd t}\int_V u\,\dd V
    =-\int_{\partial V}\qvec\cdot\nvec\,\dd A
    +\int_V \dot{q}_e\,\dd V,
\end{equation}
where \(\dot{q}_e\) is a volumetric heat source. The first term on the right is energy leaving the volume by heat flow. Using the divergence theorem,
\begin{equation}
    \int_{\partial V}\qvec\cdot\nvec\,\dd A
    =\int_V \nabla\cdot\qvec\,\dd V.
\end{equation}
Because the control volume is arbitrary,
\begin{equation}
    \boxed{\frac{\partial u}{\partial t}+\nabla\cdot\qvec=\dot{q}_e.}
    \label{eq:local_energy_conservation}
\end{equation}
This equation is exact at the continuum energy-balance level. The modeling question is how to compute \(\qvec\).

\section{Fourier conduction as a local constitutive model}
Fourier conduction assumes
\begin{equation}
    \boxed{\qvec=-k\nabla T.}
    \label{eq:fourier_law}
\end{equation}
The minus sign means heat flows from high temperature to low temperature. If \(u=C T\) with constant \(C\), substitution into Eq.~\eqref{eq:local_energy_conservation} gives
\begin{equation}
    C\frac{\partial T}{\partial t}+\nabla\cdot(-k\nabla T)=\dot{q}_e,
\end{equation}
so
\begin{equation}
    \boxed{C\frac{\partial T}{\partial t}=\nabla\cdot(k\nabla T)+\dot{q}_e.}
    \label{eq:fourier_heat_equation_general}
\end{equation}
If \(k\) is constant and there is no source,
\begin{equation}
    \frac{\partial T}{\partial t}=\alpha \nabla^2T,
    \qquad
    \alpha=\frac{k}{C}.
\end{equation}

\subsection{Why the Fourier heat equation has infinite propagation speed}
For an infinite one-dimensional medium, the Green's function of
\begin{equation}
    \frac{\partial T}{\partial t}=\alpha\frac{\partial^2T}{\partial x^2}
\end{equation}
from a point impulse at \(t=0\) is
\begin{equation}
    G(x,t)=\frac{1}{\sqrt{4\pi\alpha t}}\exp\left(-\frac{x^2}{4\alpha t}\right), \qquad t>0.
\end{equation}
For every \(t>0\), the exponential is nonzero for every finite \(x\). The disturbance immediately has nonzero amplitude arbitrarily far away. Mathematically, this is because the heat equation is parabolic. Physically, this is incompatible with heat being carried by microscopic carriers moving at finite group velocity.

At macroscale, this infinite-speed artifact is usually harmless. At nanoscale, microscale, or ultrafast timescales, it can dominate early-time predictions such as surface heat flux.

\section{The Knudsen number and transport regimes}
The mean free path \(\mfp\) is the average distance a carrier travels before scattering. The relaxation time \(\tau\) is the average time between scattering events. If the carrier speed is \(v\), then
\begin{equation}
    \boxed{\mfp=v\tau.}
    \label{eq:mfp_def}
\end{equation}
The Knudsen number is
\begin{equation}
    \boxed{\mathrm{Kn}=\frac{\mfp}{L}.}
    \label{eq:kn_def}
\end{equation}
It compares the microscopic distance between scattering events with the macroscopic device length.

\begin{center}
\begin{tabularx}{0.95\textwidth}{>{\RaggedRight\arraybackslash}p{0.22\textwidth}>{\RaggedRight\arraybackslash}p{0.25\textwidth}>{\RaggedRight\arraybackslash}X}
\toprule
Regime & Scaling & Physical interpretation \\
\midrule
Diffusive & \(\mathrm{Kn}\ll 1\) & Carriers scatter many times before crossing the system. Local equilibrium is a good approximation. Fourier conduction is usually valid. \\
Quasi-ballistic & \(\mathrm{Kn}\sim 0.1-1\) & A significant fraction of carriers travel across device-scale distances before scattering. Boundary memory and finite flight time matter. \\
Ballistic & \(\mathrm{Kn}\gtrsim 1\) & Carriers can cross the domain with few or no scattering events. Pure diffusion is a poor approximation. \\
\bottomrule
\end{tabularx}
\end{center}

\section{Microscopic starting point: distribution function}
Let \(f(t,\rvec,\bm{k})\) be the carrier distribution function, where \(\bm{k}\) is wave vector. For a phonon branch, the angular frequency is \(\omega=\omega(\bm{k})\), and the group velocity is
\begin{equation}
    \bm{v}=\nabla_{\bm{k}}\omega(\bm{k}).
\end{equation}
The quantity \(f\) is an occupation number: it tells us how many carriers occupy a phase-space state. For phonons at local equilibrium,
\begin{equation}
    f_0(\omega,T)=\frac{1}{\exp(\hbar\omega/\kB T)-1},
\end{equation}
which is the Bose-Einstein distribution. For electrons, the equilibrium distribution would be Fermi-Dirac. Module 1 focuses on the phonon-like heat-carrier derivation, but the transport structure is more general.

\section{Deriving the Boltzmann transport equation}
\subsection{Streaming without scattering}
First ignore collisions. If carriers move with velocity \(\bm{v}\), then the distribution is conserved along a trajectory in phase space. In the absence of external force, the wave vector is constant and
\begin{equation}
    \frac{\dd f}{\dd t}
    =\frac{\partial f}{\partial t}
    +\frac{\dd \rvec}{\dd t}\cdot\nabla_{\rvec}f
    =\frac{\partial f}{\partial t}+\bm{v}\cdot\nabla_{\rvec}f.
\end{equation}
Conservation along trajectories gives
\begin{equation}
    \frac{\partial f}{\partial t}+\bm{v}\cdot\nabla_{\rvec}f=0.
\end{equation}
This is free streaming.

\subsection{Adding scattering}
Scattering changes the occupation of a state. The general Boltzmann equation is
\begin{equation}
    \frac{\partial f}{\partial t}+\bm{v}\cdot\nabla_{\rvec}f
    =\left(\frac{\partial f}{\partial t}\right)_{\mathrm{coll}},
\end{equation}
where the right side is the collision operator. The exact collision operator can be complicated because it can include phonon-phonon scattering, phonon-boundary scattering, impurity scattering, electron-phonon scattering, and other mechanisms.

The relaxation-time approximation replaces the detailed collision operator by a tendency to relax toward local equilibrium:
\begin{equation}
    \left(\frac{\partial f}{\partial t}\right)_{\mathrm{coll}}
    =-\frac{f-f_0}{\tau_\omega}.
\end{equation}
Therefore,
\begin{equation}
    \boxed{\frac{\partial f}{\partial t}+\bm{v}\cdot\nabla_{\rvec}f
    =-\frac{f-f_0}{\tau_\omega}.}
    \label{eq:bte_rta}
\end{equation}
This is the Boltzmann transport equation under the relaxation-time approximation.

\subsection{Term-by-term meaning}
Each term in Eq.~\eqref{eq:bte_rta} has a direct physical meaning:
\begin{align}
    \frac{\partial f}{\partial t} &: \text{local time change of carrier occupation,}\\
    \bm{v}\cdot\nabla_{\rvec}f &: \text{change due to carriers streaming through space,}\\
    -\frac{f-f_0}{\tau_\omega} &: \text{relaxation toward local equilibrium by scattering.}
\end{align}
If \(f>f_0\), scattering reduces \(f\). If \(f<f_0\), scattering increases \(f\). The relaxation time \(\tau_\omega\) can depend on frequency because different phonon modes scatter differently.

\begin{notebox}
\textbf{Assumption in this module.} We neglect explicit force terms. For electron transport in an electric field, the BTE would include a force term in wave-vector space. For phonon heat transport without external body force, the streaming-plus-collision form is the relevant starting point.
\end{notebox}

\section{Why a direct Boltzmann solve is expensive}
In three-dimensional space, \(f\) depends on
\begin{equation}
    (x,y,z),\quad \text{two direction angles},\quad \omega,\quad t.
\end{equation}
Equivalently, the BTE lives in real space, momentum space, and time. Solving it directly is much more expensive than solving a heat equation for \(T(x,y,z,t)\). The purpose of the ballistic-diffusive approximation is to keep the finite-speed and nonlocal physics that matter in small structures while reducing the problem to a PDE in ordinary space and time plus an explicitly computed ballistic flux term.

\section{Intensity notation}
\subsection{Definition of spectral intensity}
It is useful to rewrite the BTE in terms of spectral intensity. For an isotropic medium, define
\begin{equation}
    I_\omega(t,\rvec,\hat{\Omega})
    =\frac{1}{4\pi}\,v\,\hbar\omega\,D(\omega)\,f(t,\rvec,\hat{\Omega},\omega),
    \label{eq:intensity_definition}
\end{equation}
where \(D(\omega)\) is the density of states per unit volume. The factor \(1/(4\pi)\) distributes the isotropic density over solid angle. The exact convention can vary, but this convention makes the angular moments transparent.

The intensity form of the BTE is
\begin{equation}
    \boxed{\frac{\partial I_\omega}{\partial t}+\bm{v}\cdot\nabla I_\omega
    =-\frac{I_\omega-I_{0\omega}}{\tau_\omega}.}
    \label{eq:intensity_bte}
\end{equation}
Here \(I_{0\omega}\) is the local-equilibrium intensity.

\subsection{Moments of intensity}
The internal energy density is obtained by integrating energy per unit volume over frequency and direction. Because intensity is energy flow per area per time per solid angle per frequency, division by speed converts flux-like intensity into energy density:
\begin{equation}
    \boxed{u(t,\rvec)=\int_0^\infty\int_{4\pi}\frac{I_\omega(t,\rvec,\hat{\Omega})}{v}\,\dd\Omega\,\dd\omega.}
    \label{eq:energy_moment}
\end{equation}
The heat flux is the directional projection of energy flow:
\begin{equation}
    \boxed{\qvec(t,\rvec)=\int_0^\infty\int_{4\pi} I_\omega(t,\rvec,\hat{\Omega})\,\hat{\Omega}\,\dd\Omega\,\dd\omega.}
    \label{eq:flux_moment}
\end{equation}
These two moment definitions are the bridge between microscopic carrier transport and continuum thermal variables.

\section{Recovering Fourier conduction from the BTE}
This section is important because it shows what Fourier conduction throws away. Suppose the system is close to local equilibrium, so
\begin{equation}
    f=f_0+f_1,
    \qquad
    \abs{f_1}\ll\abs{f_0}.
\end{equation}
Assume steady or slowly varying conditions so the leading nonequilibrium correction comes from spatial gradients. The BTE gives approximately
\begin{equation}
    \bm{v}\cdot\nabla f_0\approx -\frac{f_1}{\tau_\omega},
\end{equation}
so
\begin{equation}
    f_1\approx -\tau_\omega\bm{v}\cdot\nabla f_0.
    \label{eq:f1_fourier}
\end{equation}
Because \(f_0=f_0[T(\rvec)]\),
\begin{equation}
    \nabla f_0=\frac{\partial f_0}{\partial T}\nabla T.
\end{equation}
Thus
\begin{equation}
    f_1\approx -\tau_\omega(\bm{v}\cdot\nabla T)\frac{\partial f_0}{\partial T}.
\end{equation}
The heat flux is the first velocity moment of the nonequilibrium correction. In an isotropic medium,
\begin{equation}
    \int v_i v_j\,\dd\Omega = \frac{4\pi}{3}v^2\delta_{ij}.
\end{equation}
Therefore the flux becomes
\begin{equation}
    \qvec=-\frac{1}{3}\int_0^\infty C_\omega v^2\tau_\omega\,\dd\omega\,\nabla T.
\end{equation}
Using \(\mfp_\omega=v\tau_\omega\), this is
\begin{equation}
    \boxed{\qvec=-k\nabla T,\qquad
    k=\frac{1}{3}\int_0^\infty C_\omega v\mfp_\omega\,\dd\omega.}
    \label{eq:kinetic_k}
\end{equation}
This is Fourier's law from kinetic theory.

\begin{physicsbox}
\textbf{What had to be assumed to get Fourier's law?} We assumed near local equilibrium, small anisotropy, many scattering events over the temperature-gradient length scale, and no strong boundary memory. Those assumptions are precisely what become questionable when \(\mathrm{Kn}\) is not small.
\end{physicsbox}

\section{Cattaneo's correction and its limitation}
\subsection{Flux relaxation}
The Cattaneo model modifies Fourier's law by giving heat flux a relaxation time:
\begin{equation}
    \boxed{\tau\frac{\partial \qvec}{\partial t}+\qvec=-k\nabla T.}
    \label{eq:cattaneo_law}
\end{equation}
If \(\tau\to 0\), Fourier's law is recovered. Combining Eq.~\eqref{eq:cattaneo_law} with energy conservation gives a hyperbolic heat equation. For constant \(C\), constant \(k\), and no source,
\begin{equation}
    C\frac{\partial T}{\partial t}+\nabla\cdot\qvec=0.
\end{equation}
Take a time derivative:
\begin{equation}
    C\frac{\partial^2T}{\partial t^2}+\nabla\cdot\frac{\partial\qvec}{\partial t}=0.
\end{equation}
Take the divergence of Eq.~\eqref{eq:cattaneo_law}:
\begin{equation}
    \tau\nabla\cdot\frac{\partial\qvec}{\partial t}+\nabla\cdot\qvec=-k\nabla^2T.
\end{equation}
Using \(\nabla\cdot\qvec=-C\partial T/\partial t\), we get
\begin{equation}
    \boxed{\tau\frac{\partial^2T}{\partial t^2}+\frac{\partial T}{\partial t}=\alpha\nabla^2T.}
    \label{eq:cattaneo_heat_equation}
\end{equation}
This is a telegraph-type equation. It has finite signal speed
\begin{equation}
    c_{\mathrm{th}}=\sqrt{\frac{\alpha}{\tau}}.
\end{equation}
Using \(\alpha=k/C\) and the kinetic estimate \(k\approx C v\mfp/3=Cv^2\tau/3\),
\begin{equation}
    c_{\mathrm{th}}\approx \frac{v}{\sqrt{3}}.
\end{equation}
So Cattaneo fixes the most obvious mathematical pathology of Fourier conduction: infinite propagation speed.

\subsection{Why Cattaneo is still not enough}
Cattaneo still treats the entire heat flux as a local constitutive variable that relaxes toward \(-k\nabla T\). It does not explicitly represent carriers that left a boundary, crossed the domain, and arrived at an interior point without scattering. Therefore, it does not contain the correct nonlocal boundary memory.

This missing ingredient matters in thin films and nanoscale structures. In a heating experiment where one boundary suddenly emits carriers, some carriers travel ballistically into the film before enough scattering has occurred to establish a local gradient-driven flux. Cattaneo can produce wave-like behavior, but that wave is not the same as a distribution of particles arriving from the boundary at retarded times. In Chen's comparison calculations, Cattaneo gives artificial surface heat-flux oscillations and poor agreement with the Boltzmann equation in the small-scale transient regime.

\section{Ballistic-diffusive idea}
The ballistic-diffusive model starts from a simple but powerful split:
\begin{equation}
    \boxed{I_\omega=I_{b\omega}+I_{m\omega}.}
    \label{eq:bd_intensity_split}
\end{equation}
The two parts have different physical meanings:
\begin{align}
    I_{b\omega} &: \text{carriers that originated at a boundary and have not yet scattered,}\\
    I_{m\omega} &: \text{carriers produced by scattering or emission inside the medium.}
\end{align}

\begin{physicsbox}
\textbf{Why the split is useful.} The ballistic part is highly directional and nonlocal, so we should not approximate it by diffusion. The medium part has already undergone scattering and is more nearly isotropic, so a diffusion-like angular closure is more defensible.
\end{physicsbox}

This is not merely a mathematical trick. It corresponds to two populations of heat carriers with different histories:
\begin{enumerate}[leftmargin=1.6em]
    \item A boundary-originating population that remembers the boundary condition and flight time.
    \item An internally scattered population that has partially lost directional memory.
\end{enumerate}

\section{Equation for the ballistic component}
The ballistic component experiences streaming and out-scattering, but it has no in-scattering source because newly scattered carriers are assigned to \(I_{m\omega}\). Therefore,
\begin{equation}
    \boxed{\frac{1}{v}\frac{\partial I_{b\omega}}{\partial t}
    +\hat{\Omega}\cdot\nabla I_{b\omega}
    =-\frac{I_{b\omega}}{v\tau_\omega}
    =-\frac{I_{b\omega}}{\mfp_\omega}.}
    \label{eq:ballistic_eq}
\end{equation}

\subsection{Derivation along a characteristic}
Let \(s\) be distance along a ray in direction \(\hat{\Omega}\). Along the ray,
\begin{equation}
    \rvec(s)=\rvec_0+s\hat{\Omega}.
\end{equation}
For a time-dependent problem, a carrier travels a distance \(\dd s\) in time \(\dd t=\dd s/v\). Along the characteristic,
\begin{equation}
    \frac{\dd I_{b\omega}}{\dd s}
    =\frac{1}{v}\frac{\partial I_{b\omega}}{\partial t}
    +\hat{\Omega}\cdot\nabla I_{b\omega}.
\end{equation}
Therefore Eq.~\eqref{eq:ballistic_eq} becomes
\begin{equation}
    \frac{\dd I_{b\omega}}{\dd s}=-\frac{I_{b\omega}}{\mfp_\omega}.
\end{equation}
The solution is
\begin{equation}
    I_{b\omega}(s)=I_{b\omega}(s_0)\exp\left[-\int_{s_0}^{s}\frac{\dd s'}{\mfp_\omega(s')}\right].
\end{equation}
For a homogeneous medium, this reduces to
\begin{equation}
    I_{b\omega}(s)=I_{b\omega}(s_0)\exp\left[-\frac{s-s_0}{\mfp_\omega}\right].
\end{equation}
The exponential is a survival probability: it is the probability that a carrier travels from \(s_0\) to \(s\) without scattering.

\subsection{Retarded boundary time}
If the boundary intensity changes with time, the carrier arriving at \((t,\rvec)\) did not leave the boundary at time \(t\). It left earlier by the flight time
\begin{equation}
    \Delta t_{\mathrm{flight}}=\frac{s-s_0}{v}.
\end{equation}
Thus the ballistic solution is
\begin{equation}
    \boxed{I_{b\omega}(t,\rvec,\hat{\Omega})
    =I_{w\omega}\left(t-\frac{s-s_0}{v},\rvec-(s-s_0)\hat{\Omega}\right)
    \exp\left[-\int_{s_0}^{s}\frac{\dd s'}{\mfp_\omega(s')}\right].}
    \label{eq:ballistic_solution}
\end{equation}
This is one of the most important equations in the module. It contains both finite propagation speed and attenuation by scattering.

\section{Equation for the medium/scattered component}
The remaining carriers are assigned to the medium component. They obey
\begin{equation}
    \boxed{\frac{\partial I_{m\omega}}{\partial t}
    +\bm{v}\cdot\nabla I_{m\omega}
    =-\frac{I_{m\omega}-I_{0\omega}}{\tau_\omega}.}
    \label{eq:medium_eq}
\end{equation}
This equation still resembles the BTE, but now the distribution is more isotropic than the full distribution because the strongly directional unscattered boundary population has been separated out.

\section{P1 angular expansion for the medium component}
The diffusion approximation is an angular-moment closure. We expand \(I_{m\omega}\) in angular harmonics and keep only the isotropic term and the first directional term:
\begin{equation}
    \boxed{I_{m\omega}(t,\rvec,\hat{\Omega})
    \approx J_{0\omega}(t,\rvec)+\bm{J}_{1\omega}(t,\rvec)\cdot\hat{\Omega}.}
    \label{eq:p1_expansion}
\end{equation}
The scalar \(J_{0\omega}\) represents the angular average. The vector \(\bm{J}_{1\omega}\) represents the leading directional bias.

The angular identities needed are
\begin{align}
    \int_{4\pi}\dd\Omega &=4\pi,\\
    \int_{4\pi}\hat{\Omega}\,\dd\Omega &=\bm{0},\\
    \int_{4\pi}\hat{\Omega}_i\hat{\Omega}_j\,\dd\Omega &=\frac{4\pi}{3}\delta_{ij}.
    \label{eq:angular_identities}
\end{align}
These identities encode isotropy of the sphere. The factor \(1/3\) is the same factor that appears in kinetic theory expressions such as \(k=(1/3)Cv\mfp\).

\section{Moment equations for the medium component}
Substitute Eq.~\eqref{eq:p1_expansion} into Eq.~\eqref{eq:medium_eq}. Since \(\bm{v}=v\hat{\Omega}\),
\begin{equation}
    \frac{\partial J_{0\omega}}{\partial t}
    +\frac{\partial \bm{J}_{1\omega}}{\partial t}\cdot\hat{\Omega}
    +v\hat{\Omega}\cdot\nabla J_{0\omega}
    +v\hat{\Omega}\cdot\nabla(\bm{J}_{1\omega}\cdot\hat{\Omega})
    =-\frac{J_{0\omega}-I_{0\omega}}{\tau_\omega}
    -\frac{\bm{J}_{1\omega}\cdot\hat{\Omega}}{\tau_\omega}.
    \label{eq:substituted_p1}
\end{equation}

\subsection{Zeroth angular moment}
Integrate Eq.~\eqref{eq:substituted_p1} over all directions. Terms odd in \(\hat{\Omega}\) vanish. Using Eq.~\eqref{eq:angular_identities},
\begin{equation}
    4\pi\frac{\partial J_{0\omega}}{\partial t}
    +\frac{4\pi}{3}v\nabla\cdot\bm{J}_{1\omega}
    =-\frac{4\pi}{\tau_\omega}(J_{0\omega}-I_{0\omega}).
    \label{eq:zeroth_moment_dimensional}
\end{equation}
Equivalently,
\begin{equation}
    \frac{1}{v}\frac{\partial J_{0\omega}}{\partial t}
    +\frac{1}{3}\nabla\cdot\bm{J}_{1\omega}
    =-\frac{J_{0\omega}-I_{0\omega}}{v\tau_\omega}
    =-\frac{J_{0\omega}-I_{0\omega}}{\mfp_\omega}.
\end{equation}

\subsection{First angular moment}
Multiply Eq.~\eqref{eq:substituted_p1} by \(\hat{\Omega}\) and integrate over all directions. The result is
\begin{equation}
    \frac{1}{v}\frac{\partial \bm{J}_{1\omega}}{\partial t}+\nabla J_{0\omega}
    =-\frac{\bm{J}_{1\omega}}{\mfp_\omega}.
    \label{eq:first_moment_p1}
\end{equation}
This equation says that the directional bias \(\bm{J}_{1\omega}\) relaxes over a mean free path and is driven by gradients of the angular average \(J_{0\omega}\).

\section{Heat flux and internal energy split}
Using the intensity moments, the total internal energy density is
\begin{equation}
    u=u_b+u_m,
    \label{eq:u_split}
\end{equation}
where
\begin{align}
    u_b &=\int_0^\infty\int_{4\pi}\frac{I_{b\omega}}{v}\,\dd\Omega\,\dd\omega,\\
    u_m &=\int_0^\infty\int_{4\pi}\frac{I_{m\omega}}{v}\,\dd\Omega\,\dd\omega.
\end{align}
With the P1 expansion,
\begin{equation}
    u_m=\int_0^\infty \frac{4\pi J_{0\omega}}{v}\,\dd\omega.
\end{equation}
Similarly, the total heat flux is
\begin{equation}
    \qvec=\qvec_b+\qvec_m,
    \label{eq:q_split}
\end{equation}
where
\begin{align}
    \qvec_b &=\int_0^\infty\int_{4\pi} I_{b\omega}\hat{\Omega}\,\dd\Omega\,\dd\omega,\\
    \qvec_m &=\int_0^\infty\int_{4\pi} I_{m\omega}\hat{\Omega}\,\dd\Omega\,\dd\omega.
\end{align}
Using the P1 expansion and angular identities,
\begin{equation}
    \qvec_m=\frac{4\pi}{3}\int_0^\infty \bm{J}_{1\omega}\,\dd\omega.
    \label{eq:qm_j1}
\end{equation}

\section{Constitutive relation for the medium component}
Start from the first moment equation,
\begin{equation}
    \frac{1}{v}\frac{\partial \bm{J}_{1\omega}}{\partial t}+\nabla J_{0\omega}
    =-\frac{\bm{J}_{1\omega}}{\mfp_\omega}.
\end{equation}
Multiply by the spectral weight that converts \(\bm{J}_{1\omega}\) into heat flux and integrate over frequency. Under the standard averaged-parameter approximation, the result is
\begin{equation}
    \boxed{\tau\frac{\partial \qvec_m}{\partial t}+\qvec_m
    =-\frac{k}{C}\nabla u_m.}
    \label{eq:qm_constitutive_um}
\end{equation}
If we define \(T_m=u_m/C\), then
\begin{equation}
    \boxed{\tau\frac{\partial \qvec_m}{\partial t}+\qvec_m=-k\nabla T_m.}
    \label{eq:qm_constitutive_Tm}
\end{equation}

This looks like Cattaneo's law, but the interpretation is different. In Cattaneo theory, \(\qvec\) is the total heat flux. Here, \(\qvec_m\) is only the scattered medium flux. The ballistic flux \(\qvec_b\) is not represented by \(-k\nabla T_m\); it is computed separately from the boundary-originating carrier population.

\section{Energy equation for the medium component}
Use total energy conservation,
\begin{equation}
    \frac{\partial (u_m+u_b)}{\partial t}+\nabla\cdot(\qvec_m+\qvec_b)=\dot{q}_e.
\end{equation}
Rearrange to isolate the medium component:
\begin{equation}
    \frac{\partial u_m}{\partial t}+\nabla\cdot\qvec_m
    =\dot{q}_e-\frac{\partial u_b}{\partial t}-\nabla\cdot\qvec_b.
    \label{eq:medium_energy_balance}
\end{equation}
The right side says that the medium energy changes not only due to the external source, but also due to energy removed from or added by the ballistic component.

\subsection{Ballistic moment relation}
The ballistic intensity equation implies a moment equation. Integrating Eq.~\eqref{eq:ballistic_eq} over direction and frequency gives
\begin{equation}
    \frac{\partial u_b}{\partial t}+\nabla\cdot\qvec_b=-\frac{u_b}{\tau}
    \label{eq:ballistic_moment_unmultiplied}
\end{equation}
for an averaged relaxation-time model. Equivalently,
\begin{equation}
    \tau\frac{\partial u_b}{\partial t}+\tau\nabla\cdot\qvec_b=-u_b.
    \label{eq:ballistic_moment}
\end{equation}
This states that the ballistic energy population is depleted by scattering over time \(\tau\). Once a ballistic carrier scatters, it leaves the ballistic population and becomes part of the medium population.

\subsection{Eliminating the medium heat flux}
Take the divergence of Eq.~\eqref{eq:qm_constitutive_um}:
\begin{equation}
    \tau\frac{\partial}{\partial t}(\nabla\cdot\qvec_m)+\nabla\cdot\qvec_m
    =-\nabla\cdot\left(\frac{k}{C}\nabla u_m\right).
    \label{eq:div_qm_constitutive}
\end{equation}
From Eq.~\eqref{eq:medium_energy_balance},
\begin{equation}
    \nabla\cdot\qvec_m=\dot{q}_e-\frac{\partial u_m}{\partial t}-\frac{\partial u_b}{\partial t}-\nabla\cdot\qvec_b.
    \label{eq:div_qm_from_energy}
\end{equation}
Substitute Eq.~\eqref{eq:div_qm_from_energy} into Eq.~\eqref{eq:div_qm_constitutive}. After collecting terms,
\begin{align}
    \tau\frac{\partial^2u_m}{\partial t^2}+\frac{\partial u_m}{\partial t}
    &=\nabla\cdot\left(\frac{k}{C}\nabla u_m\right)
    +\dot{q}_e+\tau\frac{\partial\dot{q}_e}{\partial t}
    -\nabla\cdot\qvec_b \\
    &\quad -\left[\tau\frac{\partial^2u_b}{\partial t^2}+\frac{\partial u_b}{\partial t}
    +\tau\frac{\partial}{\partial t}(\nabla\cdot\qvec_b)\right].
\end{align}
The bracketed term vanishes by differentiating Eq.~\eqref{eq:ballistic_moment}. Therefore,
\begin{equation}
    \boxed{\tau\frac{\partial^2u_m}{\partial t^2}+\frac{\partial u_m}{\partial t}
    =\nabla\cdot\left(\frac{k}{C}\nabla u_m\right)
    -\nabla\cdot\qvec_b
    +\left(\dot{q}_e+\tau\frac{\partial\dot{q}_e}{\partial t}\right).}
    \label{eq:bd_energy_final}
\end{equation}
This is the ballistic-diffusive heat-conduction equation in energy-density form.

\begin{physicsbox}
\textbf{Interpretation of the final equation.} The left side is a Cattaneo-like inertial thermal response of the scattered medium energy. The first term on the right is diffusion of the scattered medium energy. The second term, \(-\nabla\cdot\qvec_b\), is the explicit source/sink caused by divergence of the boundary-originating ballistic heat flux. The final terms represent external volumetric heating and its relaxation-time correction.
\end{physicsbox}

\section{Explicit ballistic heat flux}
The ballistic flux is obtained by inserting the characteristic solution Eq.~\eqref{eq:ballistic_solution} into the flux moment:
\begin{equation}
    \boxed{\qvec_b(t,\rvec)=\int_0^\infty\int_{4\pi}
    I_{w\omega}\left(t-\frac{s-s_0}{v},\rvec-(s-s_0)\hat{\Omega}\right)
    \exp\left[-\int_{s_0}^{s}\frac{\dd s'}{\mfp_\omega(s')}\right]
    \hat{\Omega}\,\dd\Omega\,\dd\omega.}
    \label{eq:ballistic_flux_general}
\end{equation}
This term depends on boundary emission/reflection, flight time, geometry, and attenuation. In simple geometries, it can be evaluated analytically or semi-analytically. In complex 3D geometries, it may be approximated by ray tracing, angular quadrature, or reduced boundary-source models.

\section{Boundary condition for the medium component}
Because boundary-originating carriers are assigned to the ballistic part, the medium component at a boundary consists only of carriers incident on the boundary from inside the material. This leads to a Marshak-type boundary condition.

At a boundary with outward normal \(\nvec\), the incident medium flux through the boundary is determined by directions satisfying \(\hat{\Omega}\cdot\nvec<0\). The P1 approximation gives
\begin{equation}
    \qvec_m\cdot\nvec=-\frac{v}{2}u_m
    \label{eq:marshak_flux_relation}
\end{equation}
for a black-wall limiting form. Combining this with the medium constitutive relation gives
\begin{equation}
    \boxed{\tau\frac{\partial u_m}{\partial t}+u_m
    =-\frac{\mfp}{3}\nabla u_m\cdot\nvec.}
    \label{eq:marshak_um}
\end{equation}
This is not the same as prescribing the boundary temperature in Fourier conduction. It is a boundary condition for the scattered medium energy after boundary-originating carriers have been assigned to the ballistic component.

\section{Temperature consistency at boundaries}
The boundary temperature that characterizes emitted carriers is not always identical to the equivalent equilibrium temperature just inside the medium. This distinction is central in nanoscale transport.

Suppose a boundary emits phonons as if it were at temperature \(T_1\). The incoming carrier population is then specified by an emitted distribution. Inside the film, however, the total local energy is
\begin{equation}
    u(x,t)=u_m(x,t)+u_b(x,t).
\end{equation}
An equivalent internal-energy temperature would be
\begin{equation}
    T_{\mathrm{eq}}(x,t)=T_0+\frac{u(x,t)-u_0}{C}
\end{equation}
for a constant heat capacity model. In a highly nonequilibrium boundary layer, \(T_{\mathrm{eq}}\) need not equal the emitted boundary temperature. Therefore, a naively imposed Dirichlet temperature boundary condition can be physically inconsistent.

This is why the ballistic-diffusive formulation keeps the boundary-emitted ballistic population separate from the internal medium population. It avoids forcing all boundary physics into a local equilibrium temperature.


\section{One-dimensional thin-film validation problem}
\subsection{Why this validation problem is the right first test}
The canonical validation problem is transient heat conduction across a film of thickness \(L\). The coordinate is \(x\in[0,L]\). The film starts at uniform reference temperature \(T_0\). At \(t=0\), the left boundary begins emitting carriers characterized by a higher emitted temperature \(T_1\), while the right boundary remains characterized by \(T_0\). There is no volumetric heating:
\begin{equation}
    \dot{q}_e=0.
\end{equation}

This problem is simple geometrically but difficult physically. It contains all of the effects that distinguish Fourier, Cattaneo, Boltzmann, and ballistic-diffusive descriptions:
\begin{enumerate}[leftmargin=1.6em]
    \item a finite device length \(L\),
    \item a finite carrier mean free path \(\mfp\),
    \item a finite carrier flight time across the film,
    \item a boundary-emitted carrier population,
    \item a scattered carrier population inside the film,
    \item a temperature or energy jump issue at the boundary.
\end{enumerate}

\begin{notebox}
\textbf{Slide takeaway.} The thin-film problem is the cleanest benchmark because it is one-dimensional in space but still contains the real nanoscale physics: finite mean free path, finite carrier speed, retarded boundary emission, and boundary temperature subtlety.
\end{notebox}

\subsection{Starting from the full ballistic-diffusive energy equation}
The general ballistic-diffusive medium-energy equation derived above is
\begin{equation}
    \tau\frac{\partial^2u_m}{\partial t^2}
    +\frac{\partial u_m}{\partial t}
    =\nabla\cdot\left(\frac{k}{C}\nabla u_m\right)
    -\nabla\cdot\qvec_b
    +\left(\dot{q}_e+\tau\frac{\partial \dot{q}_e}{\partial t}\right).
    \label{eq:thinfilm_general_start}
\end{equation}
The unknown in this equation is not the total internal energy \(u\), but the medium or scattered-carrier energy \(u_m\). The ballistic energy \(u_b\) and ballistic flux \(\qvec_b\) are computed from boundary-originating carriers. The total energy is reconstructed afterward:
\begin{equation}
    u(x,t)=u_m(x,t)+u_b(x,t).
\end{equation}

For the one-dimensional thin film, all spatial variation is along \(x\), and the ballistic flux has only an \(x\)-component:
\begin{equation}
    \qvec_b(x,t)=q_b(x,t)\,\hat{x}.
\end{equation}
Therefore,
\begin{equation}
    \nabla\cdot\qvec_b=\frac{\partial q_b}{\partial x},
    \qquad
    \nabla\cdot\left(\frac{k}{C}\nabla u_m\right)
    =\frac{\partial}{\partial x}\left(\frac{k}{C}\frac{\partial u_m}{\partial x}\right).
\end{equation}
If \(k\) and \(C\) are treated as constants, Eq.~\eqref{eq:thinfilm_general_start} reduces to
\begin{equation}
    \boxed{\tau\frac{\partial^2u_m}{\partial t^2}+\frac{\partial u_m}{\partial t}
    =\frac{k}{C}\frac{\partial^2u_m}{\partial x^2}-\frac{\partial q_b}{\partial x}.}
    \label{eq:bd_1d_dimensional_expanded}
\end{equation}

This equation has two conceptually different right-hand-side terms:
\begin{align}
    \frac{k}{C}\frac{\partial^2u_m}{\partial x^2}
    &\quad \text{is local diffusion of the scattered carrier energy,}\\
    -\frac{\partial q_b}{\partial x}
    &\quad \text{is local gain or loss caused by the divergence of ballistic flux.}
\end{align}
The ballistic term is not an optional correction. It is the mechanism by which boundary-emitted carriers deposit energy into the interior after streaming and scattering.

\subsection{Initial and boundary data in physical variables}
At \(t=0\), the film is at equilibrium temperature \(T_0\). In a constant-\(C\) model, the reference total internal energy is
\begin{equation}
    u_0=C T_0,
\end{equation}
up to an arbitrary additive reference. The physically relevant perturbation is measured relative to this state. The initial conditions are
\begin{equation}
    u(x,0)=u_0,
    \qquad
    \left.\frac{\partial u}{\partial t}\right|_{t=0}=0.
\end{equation}
Because the model separates total energy into ballistic and medium pieces, these conditions are implemented by subtracting equilibrium contributions:
\begin{equation}
    u_m(x,0)-u_{m0}(x)=0,
    \qquad
    u_b(x,0)-u_{b0}(x)=0.
\end{equation}

The left boundary condition is not ``\(T(0,t)=T_1\)'' in the Fourier sense. Instead, the left boundary emits a population of carriers characterized by \(T_1\). The right boundary emits a population characterized by \(T_0\). The medium component satisfies the Marshak-type boundary condition derived from the one-sided incoming angular population:
\begin{align}
    \left(\tau\frac{\partial u_m}{\partial t}+u_m\right)_{x=0}
    &= -\frac{\mfp}{3}\left.\frac{\partial u_m}{\partial x}\right|_{x=0},\label{eq:left_marshak_dim}\\
    \left(\tau\frac{\partial u_m}{\partial t}+u_m\right)_{x=L}
    &= +\frac{\mfp}{3}\left.\frac{\partial u_m}{\partial x}\right|_{x=L},\label{eq:right_marshak_dim}
\end{align}
where the sign changes because the outward normal is \(-\hat{x}\) at \(x=0\) and \(+\hat{x}\) at \(x=L\). Some papers write these signs using a local outward normal convention. The safest implementation rule is to evaluate the normal derivative \(\nabla u_m\cdot\nvec\) at each boundary.

\begin{physicsbox}
\textbf{Physical meaning of the boundary condition.} The boundary condition does not say that the medium energy is fixed at the wall. It says that the diffusive medium flux at a boundary is made only of the carriers incident on that boundary from inside the domain. Carriers emitted by the wall are assigned to the ballistic component.
\end{physicsbox}

\subsection{Nondimensional variables}
The thin-film equation becomes much more interpretable after nondimensionalization. Define
\begin{equation}
    \eta=\frac{x}{L},
    \qquad
    t^*=\frac{t}{\tau},
    \qquad
    \mathrm{Kn}=\frac{\mfp}{L},
    \qquad
    \Delta T=T_1-T_0.
\end{equation}
Here \(\eta\) measures position as a fraction of the film thickness, \(t^*\) measures time in units of one scattering time, and \(\mathrm{Kn}\) measures how far a carrier travels before scattering compared with the film thickness.

The normalized medium energy is
\begin{equation}
    \theta_m(\eta,t^*)=\frac{u_m(x,t)-u_{m0}(x)}{C\Delta T}.
    \label{eq:theta_m_def_expanded}
\end{equation}
The normalized ballistic energy and total energy are
\begin{equation}
    \theta_b(\eta,t^*)=\frac{u_b(x,t)-u_{b0}(x)}{C\Delta T},
    \qquad
    \theta(\eta,t^*)=\theta_m(\eta,t^*)+\theta_b(\eta,t^*).
\end{equation}
The normalized ballistic flux is
\begin{equation}
    q_b^*(\eta,t^*)=\frac{q_b(x,t)-q_{b0}(x)}{C v\Delta T}.
    \label{eq:qb_star_def_expanded}
\end{equation}

The derivative conversions are
\begin{equation}
    \frac{\partial}{\partial t}=\frac{1}{\tau}\frac{\partial}{\partial t^*},
    \qquad
    \frac{\partial^2}{\partial t^2}=\frac{1}{\tau^2}\frac{\partial^2}{\partial {t^*}^2},
\end{equation}
\begin{equation}
    \frac{\partial}{\partial x}=\frac{1}{L}\frac{\partial}{\partial\eta},
    \qquad
    \frac{\partial^2}{\partial x^2}=\frac{1}{L^2}\frac{\partial^2}{\partial\eta^2}.
\end{equation}

We also use the kinetic-theory estimate for thermal conductivity:
\begin{equation}
    k \approx \frac{1}{3}C v\mfp.
    \label{eq:kinetic_k_thinfilm}
\end{equation}
The factor \(1/3\) appears because in an isotropic three-dimensional carrier gas, only one third of the squared velocity contributes to transport along a chosen coordinate direction:
\begin{equation}
    \left< v_x^2\right>=\frac{1}{3}v^2.
\end{equation}
Since \(\mfp=v\tau\), Eq.~\eqref{eq:kinetic_k_thinfilm} is equivalently
\begin{equation}
    \frac{k}{C}=\frac{v\mfp}{3}=\frac{v^2\tau}{3}.
\end{equation}

\subsection{Step-by-step nondimensionalization of the governing equation}
Start with
\begin{equation}
    \tau\frac{\partial^2u_m}{\partial t^2}+\frac{\partial u_m}{\partial t}
    =\frac{k}{C}\frac{\partial^2u_m}{\partial x^2}-\frac{\partial q_b}{\partial x}.
\end{equation}
Using \(u_m-u_{m0}=C\Delta T\theta_m\), the left-hand side becomes
\begin{align}
    \tau\frac{\partial^2u_m}{\partial t^2}+\frac{\partial u_m}{\partial t}
    &=\tau C\Delta T\frac{1}{\tau^2}\frac{\partial^2\theta_m}{\partial {t^*}^2}
    +C\Delta T\frac{1}{\tau}\frac{\partial\theta_m}{\partial t^*}\nonumber\\
    &=\frac{C\Delta T}{\tau}
    \left(\frac{\partial^2\theta_m}{\partial {t^*}^2}
    +\frac{\partial\theta_m}{\partial t^*}\right).
    \label{eq:lhs_nd_thinfilm}
\end{align}
The diffusive term becomes
\begin{align}
    \frac{k}{C}\frac{\partial^2u_m}{\partial x^2}
    &=\frac{k}{C}\,C\Delta T\frac{1}{L^2}\frac{\partial^2\theta_m}{\partial\eta^2}\\
    &=k\Delta T\frac{1}{L^2}\frac{\partial^2\theta_m}{\partial\eta^2}.
\end{align}
Divide this by the common left-hand-side scale \(C\Delta T/\tau\):
\begin{align}
    \frac{k\Delta T/L^2}{C\Delta T/\tau}
    &=\frac{k\tau}{C L^2}
    =\frac{1}{L^2}\left(\frac{v\mfp}{3}\right)\tau
    =\frac{v^2\tau^2}{3L^2}
    =\frac{\mfp^2}{3L^2}
    =\frac{\mathrm{Kn}^2}{3}.
\end{align}
Thus the nondimensional diffusion coefficient is \(\mathrm{Kn}^2/3\).

For the ballistic term, use \(q_b-q_{b0}=C v\Delta T\,q_b^*\):
\begin{equation}
    \frac{\partial q_b}{\partial x}
    =C v\Delta T\frac{1}{L}\frac{\partial q_b^*}{\partial\eta}.
\end{equation}
Divide by \(C\Delta T/\tau\):
\begin{equation}
    \frac{C v\Delta T/L}{C\Delta T/\tau}
    =\frac{v\tau}{L}=\frac{\mfp}{L}=\mathrm{Kn}.
\end{equation}
Therefore Eq.~\eqref{eq:bd_1d_dimensional_expanded} becomes
\begin{equation}
    \boxed{
    \frac{\partial^2\theta_m}{\partial {t^*}^2}
    +\frac{\partial\theta_m}{\partial t^*}
    =\frac{\mathrm{Kn}^2}{3}\frac{\partial^2\theta_m}{\partial\eta^2}
    -\mathrm{Kn}\frac{\partial q_b^*}{\partial\eta}.}
    \label{eq:bd_1d_dimensionless_expanded}
\end{equation}

\begin{notebox}
\textbf{Slide takeaway.} The nondimensional thin-film equation exposes the controlling parameter. The same physical equation changes character as \(\mathrm{Kn}=\mfp/L\) changes. This is why Module 1 validation must sweep \(\mathrm{Kn}\).
\end{notebox}

\subsection{Dimensionless initial and boundary conditions}
The normalized initial conditions are
\begin{equation}
    \theta(\eta,0)=0,
    \qquad
    \left.\frac{\partial\theta}{\partial t^*}\right|_{t^*=0}=0.
\end{equation}
A practical implementation usually initializes the nonequilibrium medium and ballistic perturbations as
\begin{equation}
    \theta_m(\eta,0)=0,
    \qquad
    \theta_b(\eta,0)=0,
\end{equation}
with the boundary step entering through the emitted ballistic intensity.

Using \(x=L\eta\), \(\mfp/L=\mathrm{Kn}\), and \(u_m-u_{m0}=C\Delta T\theta_m\), the Marshak-type boundary conditions become
\begin{align}
    \left(\frac{\partial\theta_m}{\partial t^*}+\theta_m\right)_{\eta=0}
    &= -\frac{\mathrm{Kn}}{3}
    \left.\frac{\partial\theta_m}{\partial\eta}\right|_{\eta=0},\label{eq:left_bc_nd}\\
    \left(\frac{\partial\theta_m}{\partial t^*}+\theta_m\right)_{\eta=1}
    &= +\frac{\mathrm{Kn}}{3}
    \left.\frac{\partial\theta_m}{\partial\eta}\right|_{\eta=1}.\label{eq:right_bc_nd}
\end{align}
Again, the sign convention is most safely implemented through the outward normal derivative.

\subsection{Ballistic intensity in one dimension}
Let \(\mu=\cos\theta\) be the directional cosine relative to the positive \(x\)-direction. A carrier emitted from the left boundary and traveling with \(0<\mu<1\) reaches location \(x\) after flight time
\begin{equation}
    t_{\mathrm{flight}}=\frac{x}{\mu v}.
\end{equation}
The path length is \(x/\mu\), not \(x\), because a shallow-angle carrier travels diagonally through the film. The probability-like attenuation factor for no scattering along that path is
\begin{equation}
    \exp\left(-\frac{x/\mu}{\mfp}\right)
    =\exp\left(-\frac{x}{\mu\mfp}\right).
\end{equation}
Therefore the left-emitted ballistic intensity has the form
\begin{equation}
    I_b^{(+)}(x,t,\mu)
    =I_{b1}\left(t-\frac{x}{\mu v}\right)
    \exp\left(-\frac{x}{\mu\mfp}\right),
    \qquad 0<\mu<1.
    \label{eq:left_ib_1d_expanded}
\end{equation}
Similarly, for right-emitted carriers moving left, \(-1<\mu<0\),
\begin{equation}
    I_b^{(-)}(x,t,\mu)
    =I_{b2}\left(t-\frac{L-x}{|\mu|v}\right)
    \exp\left(-\frac{L-x}{|\mu|\mfp}\right),
    \qquad -1<\mu<0.
    \label{eq:right_ib_1d_expanded}
\end{equation}
The time arguments are retarded boundary times. They say that the interior point sees what the boundary emitted earlier, not what the boundary emits at the current time.

\subsection{Finite arrival time and the lower angular limit}
For a step heating applied at the left boundary at \(t=0\), a left-emitted carrier contributes at position \(x\) only if
\begin{equation}
    t-\frac{x}{\mu v}\geq 0.
\end{equation}
This gives
\begin{equation}
    \mu\geq \mu_t,\qquad \mu_t=\frac{x}{vt}.
\end{equation}
In nondimensional variables, because \(x=L\eta\), \(t=\tau t^*\), and \(v\tau=\mfp=\mathrm{Kn}L\),
\begin{equation}
    \boxed{\mu_t=\frac{\eta}{\mathrm{Kn}\,t^*}.}
    \label{eq:mu_t_expanded}
\end{equation}
This expression has a direct physical interpretation:
\begin{itemize}[leftmargin=1.6em]
    \item if \(\mu_t>1\), no carrier direction has reached \(\eta\) yet;
    \item if \(0<\mu_t<1\), only sufficiently forward-directed carriers have arrived;
    \item if \(\mu_t\leq 0\), all forward directions are allowed for a sustained step source.
\end{itemize}

\subsection{Ballistic flux contribution for the step-heated film}
The heat flux is the directional energy flow. In one dimension, each carrier direction contributes a factor \(\mu\) because only the \(x\)-component of velocity transports energy across a plane normal to \(x\). For the left-emitted perturbation, the normalized ballistic flux contribution is
\begin{equation}
    \boxed{
    q_b^*(\eta,t^*)=\frac{1}{2}\int_{\mu_t}^{1}
    \mu\exp\left(-\frac{\eta}{\mu\mathrm{Kn}}\right)\,\dd\mu,}
    \qquad
    0<\mu_t<1.
    \label{eq:qb_star_integral_expanded}
\end{equation}
The factor \(1/2\) comes from angular averaging over a half-space in the one-dimensional slab reduction. When \(\mu_t>1\), the integral is zero because no ballistic carrier from the heated boundary has arrived. When \(\mu_t<0\), the lower limit is effectively zero for the forward half-space.

The derivative that enters the medium-energy equation is
\begin{equation}
    \frac{\partial q_b^*}{\partial\eta}.
\end{equation}
This term is negative where the forward ballistic flux decays with distance. Because the PDE contains \(-\mathrm{Kn}\,\partial q_b^*/\partial\eta\), attenuation of ballistic flux becomes a positive source for the scattered medium energy. In words: ballistic carriers lose energy from the unscattered population by scattering, and that energy appears in the medium population.

\subsection{Ballistic internal energy contribution}
The ballistic internal energy contribution is similar to the flux contribution but without the directional projection factor \(\mu\). A representative normalized left-boundary contribution has the form
\begin{equation}
    \theta_b(\eta,t^*)=\frac{1}{2}\int_{\mu_t}^{1}
    \exp\left(-\frac{\eta}{\mu\mathrm{Kn}}\right)\,\dd\mu,
    \qquad 0<\mu_t<1.
    \label{eq:theta_b_integral_expanded}
\end{equation}
The total normalized energy-temperature response is then
\begin{equation}
    \theta(\eta,t^*)=\theta_m(\eta,t^*)+\theta_b(\eta,t^*).
\end{equation}
This is the quantity most analogous to the temperature plotted in thin-film comparisons, with the warning that it is an equivalent internal-energy temperature in the nonequilibrium regime.

\subsection{Equilibrium subtraction and why it is not cosmetic}
The P1 approximation is an approximation to the angular distribution of the medium component. Because of this approximation, the separate equilibrium pieces \(u_{b0}\) and \(u_{m0}\) do not necessarily sum to a perfectly uniform equilibrium energy if handled naively. Chen's normalization procedure subtracts the reference equilibrium ballistic and medium components:
\begin{align}
    \theta_m &=\frac{u_m-u_{m0}}{C\Delta T},\label{eq:theta_m_subtraction}\\
    \theta_b &=\frac{u_b-u_{b0}}{C\Delta T},\label{eq:theta_b_subtraction}\\
    q_b^* &=\frac{q_b-q_{b0}}{C v\Delta T}.\label{eq:qb_subtraction}
\end{align}
This focuses the calculation on the nonequilibrium perturbation caused by the boundary heating and reduces inconsistency associated with the approximate angular closure.

One way to see the issue is to imagine the whole film at uniform equilibrium temperature \(T_0\). Physically, the total energy should be uniform. However, after splitting the angular distribution into a boundary-originating piece and an approximate P1 medium piece, the individual reconstructed pieces may have spatial dependence near boundaries. The subtraction removes the reference spatial structure and leaves only the perturbation that matters for the transient experiment.

\begin{physicsbox}
\textbf{Physical meaning.} The equilibrium subtraction is not a mathematical trick to hide an error. It is a bookkeeping correction that keeps the approximate split model focused on the nonequilibrium change from the reference state.
\end{physicsbox}

\subsection{What should be plotted in the Module 1 validation}
For \(\mathrm{Kn}=0.01,0.1,1,10\), the validation should plot:
\begin{enumerate}[leftmargin=1.6em]
    \item \(\theta_m(\eta,t^*)\): scattered medium energy,
    \item \(\theta_b(\eta,t^*)\): ballistic energy contribution,
    \item \(\theta(\eta,t^*)=\theta_m+\theta_b\): total equivalent temperature response,
    \item \(q_b^*(\eta,t^*)\): ballistic heat flux,
    \item total heat flux history at the boundaries,
    \item comparison against Fourier and Cattaneo predictions.
\end{enumerate}
The expected trend is:
\begin{align}
    \mathrm{Kn}\ll 1 &: \quad \text{Fourier-like behavior after very short transients,}\\
    \mathrm{Kn}\sim 1 &: \quad \text{strong ballistic-diffusive correction,}\\
    \mathrm{Kn}\gg 1 &: \quad \text{dominant boundary-originating ballistic transport.}
\end{align}

\section{Dedicated comparison: why Fourier and Cattaneo fail}
\subsection{Failure mode of Fourier conduction}
Fourier conduction fails in the quasi-ballistic regime for three related reasons:
\begin{enumerate}[leftmargin=1.6em]
    \item \textbf{Infinite speed.} The parabolic heat equation predicts an instantaneous response everywhere.
    \item \textbf{No carrier flight time.} It cannot represent the delay \((s-s_0)/v\) between boundary emission and interior arrival.
    \item \textbf{No nonlocal boundary memory.} The flux is controlled only by \(\nabla T\) at the same point, not by the boundary state from which unscattered carriers originated.
\end{enumerate}
As \(t\to 0^+\), Fourier theory can predict unrealistically large surface heat flux because it assumes an immediately established temperature gradient.

\subsection{Failure mode of Cattaneo conduction}
Cattaneo conduction improves Fourier theory by adding a heat-flux relaxation time, but it still fails to represent the population split. Its limitations are:
\begin{enumerate}[leftmargin=1.6em]
    \item \textbf{All flux is treated as one local variable.} There is no separate \(\qvec_b\) that remembers boundary origin.
    \item \textbf{Ballistic transport is replaced by a thermal wave.} The resulting wave speed is finite, but a thermal wave is not the same as angularly resolved particles arriving at retarded times.
    \item \textbf{Boundary conditions remain problematic.} A prescribed local temperature at a boundary is not necessarily the same as a boundary-emitted carrier distribution.
    \item \textbf{Artificial oscillations can occur.} Hyperbolic heat equations can generate nonphysical flux oscillations in regimes where the Boltzmann solution shows smooth ballistic arrival and scattering.
\end{enumerate}

\subsection{Why the ballistic-diffusive model is better}
The ballistic-diffusive model keeps the strongest advantages of both descriptions:
\begin{align}
    \text{Boltzmann-like feature:} &\quad \text{retarded boundary-originating ballistic flux } \qvec_b,\\
    \text{Diffusion-like feature:} &\quad \text{ordinary space-time PDE for the scattered medium energy } u_m.
\end{align}
It avoids a full phase-space solve while retaining the physical mechanism that Fourier and Cattaneo miss: unscattered carriers crossing a finite structure.


\section{Finite element roadmap for the ballistic-diffusive equation}
The previous sections derived the ballistic-diffusive heat-transport equation from the Boltzmann transport picture. We now turn that physics equation into the finite element method used by the code. The purpose of this section is pedagogical: we will not simply state the weak form and matrix form. We will derive the whole chain.

For the FEM derivation, it is useful to introduce a generic scalar unknown
\begin{equation}
    \phi(\rvec,t).
\end{equation}
In dimensional form, \(\phi=u_m\), the medium or scattered-carrier internal energy density. In the dimensionless one-dimensional validation problem, \(\phi=\theta_m\), the normalized medium energy. This notation lets us write one FEM derivation and then specialize it.

The full FEM chain is
\begin{equation}
\boxed{\begin{gathered}
\text{strong form}\rightarrow\text{weighted residual}\rightarrow\text{weak form}\\
\rightarrow\text{shape-function approximation}\rightarrow\text{element matrices}\\
\rightarrow\text{global assembly}\rightarrow\text{boundary conditions}\\
\rightarrow\text{nodal solve}\rightarrow\text{field reconstruction}
\end{gathered}}
\end{equation}

\begin{physicsbox}
\textbf{Core idea of FEM.} We replace the impossible task of satisfying the PDE exactly at every point by the more realistic task of satisfying an integrated, weighted form of the PDE over many small elements. The unknown field is represented by simple local interpolation functions, and the PDE becomes a matrix equation for nodal coefficients.
\end{physicsbox}

\section{The strong form that Module 1 actually solves}
\subsection{Dimensional strong form}
The dimensional ballistic-diffusive medium-energy equation is
\begin{equation}
    \tau\frac{\partial^2u_m}{\partial t^2}
    +\frac{\partial u_m}{\partial t}
    =\nabla\cdot\left(\alpha\nabla u_m\right)+S_{\mathrm{BD}},
    \qquad
    \alpha=\frac{k}{C},
    \label{eq:fem_dimensional_strong_original}
\end{equation}
where
\begin{equation}
    S_{\mathrm{BD}}
    =-\nabla\cdot\qvec_b
    +\dot q_e+\tau\frac{\partial \dot q_e}{\partial t}.
    \label{eq:fem_sbd_dimensional}
\end{equation}
Moving all differential-operator terms to the left gives the residual form
\begin{equation}
    \boxed{
    \mathcal{R}(u_m)
    =\tau\frac{\partial^2u_m}{\partial t^2}
    +\frac{\partial u_m}{\partial t}
    -\nabla\cdot(\alpha\nabla u_m)
    -S_{\mathrm{BD}}
    =0
    \quad \text{in }\Omega.}
    \label{eq:fem_dimensional_residual}
\end{equation}
This is called the \emph{strong form} because the derivatives are required pointwise. In particular, the diffusion term contains a second spatial derivative when \(\alpha\) is constant:
\begin{equation}
    \nabla\cdot(\alpha\nabla u_m)=\alpha\nabla^2u_m.
\end{equation}
A classical solution must therefore be smooth enough for these derivatives to exist pointwise. FEM weakens that requirement.

\subsection{Dimensionless one-dimensional strong form}
The first validation implementation is the nondimensional thin-film equation
\begin{equation}
    \frac{\partial^2\theta_m}{\partial {t^*}^2}
    +\frac{\partial\theta_m}{\partial t^*}
    =D\frac{\partial^2\theta_m}{\partial\eta^2}
    +s_b(\eta,t^*),
    \qquad
    D=\frac{\mathrm{Kn}^2}{3},
    \label{eq:fem_1d_bd_strong}
\end{equation}
where
\begin{equation}
    s_b(\eta,t^*)=-\mathrm{Kn}\frac{\partial q_b^*}{\partial\eta}.
    \label{eq:fem_1d_ballistic_source}
\end{equation}
Equivalently,
\begin{equation}
    \boxed{
    \mathcal{R}(\theta_m)
    =\frac{\partial^2\theta_m}{\partial {t^*}^2}
    +\frac{\partial\theta_m}{\partial t^*}
    -D\frac{\partial^2\theta_m}{\partial\eta^2}
    -s_b
    =0,
    \qquad 0<\eta<1.}
    \label{eq:fem_1d_residual}
\end{equation}
This is the simplest equation to use for teaching FEM because all spatial integrals are one-dimensional, but it already contains the important ballistic-diffusive structure: inertia-like thermal response, damping-like first time derivative, diffusion of scattered energy, and a ballistic source.

\subsection{Boundary partition}
For a general domain, the boundary is split into parts:
\begin{equation}
    \partial\Omega=\Gamma_D\cup\Gamma_N\cup\Gamma_R,
    \qquad
    \Gamma_i\cap\Gamma_j=\varnothing\quad (i\ne j).
\end{equation}
These symbols mean:
\begin{align}
    \Gamma_D &: \text{Dirichlet boundary, where the value of }\phi\text{ is prescribed},\\
    \Gamma_N &: \text{Neumann boundary, where a normal flux is prescribed},\\
    \Gamma_R &: \text{Robin or Marshak boundary, where value and normal derivative are related}.
\end{align}
For the ballistic-diffusive medium component, pure Dirichlet temperature conditions are usually not the most physical boundary model, because boundary-emitted carriers are already assigned to \(\qvec_b\). Still, Dirichlet conditions are useful for numerical tests.

\section{Why the weak form is needed}
\subsection{Pointwise residual versus weighted residual}
The strong form demands
\begin{equation}
    \mathcal{R}(\phi)=0
\end{equation}
at every point. The weighted-residual idea relaxes this. We require the residual to be orthogonal to a set of test functions \(w\):
\begin{equation}
    \int_\Omega w\,\mathcal{R}(\phi)\,\dd V=0
    \quad \text{for all admissible }w.
    \label{eq:weighted_residual_start}
\end{equation}
If this equation holds for every possible test function in a sufficiently rich space, it is equivalent to the strong form. If we restrict \(\phi\) and \(w\) to finite-dimensional approximation spaces, Eq.~\eqref{eq:weighted_residual_start} becomes a practical numerical method.

\subsection{Why integration by parts is the decisive step}
The diffusion term contains \(\nabla\cdot(\alpha\nabla\phi)\), which involves second derivatives of \(\phi\). If we used the weighted residual directly, our approximate \(\phi_h\) would need continuous second derivatives. Standard linear finite elements do not have that. They are continuous, but their slopes are discontinuous from element to element.

The key identity is the divergence product rule:
\begin{equation}
    \nabla\cdot(w\alpha\nabla\phi)
    =\nabla w\cdot\alpha\nabla\phi+w\nabla\cdot(\alpha\nabla\phi).
\end{equation}
Rearrange:
\begin{equation}
    w\nabla\cdot(\alpha\nabla\phi)
    =\nabla\cdot(w\alpha\nabla\phi)-\nabla w\cdot\alpha\nabla\phi.
\end{equation}
Integrate over \(\Omega\) and use the divergence theorem:
\begin{equation}
    \int_\Omega w\nabla\cdot(\alpha\nabla\phi)\,\dd V
    =\int_{\partial\Omega} w\alpha\nabla\phi\cdot\nvec\,\dd A
    -\int_\Omega \nabla w\cdot\alpha\nabla\phi\,\dd V.
    \label{eq:fem_parts_identity}
\end{equation}
This identity moves one spatial derivative from \(\phi\) to \(w\). After this step, the method only requires first derivatives of shape functions, which is exactly what standard FEM provides.

\begin{physicsbox}
\textbf{Physical meaning of integration by parts.} The volume diffusion operator is converted into two pieces: an interior gradient-gradient energy exchange term and a boundary normal-flux term. That boundary term is where Neumann and Marshak boundary conditions enter naturally.
\end{physicsbox}

\section{Weak form in dimensional variables}
Start from the residual equation
\begin{equation}
    \tau\phi_{tt}+\phi_t-\nabla\cdot(\alpha\nabla\phi)-S=0,
    \label{eq:fem_generic_strong}
\end{equation}
where, for the dimensional problem,
\begin{equation}
    \phi=u_m,
    \qquad
    S=S_{\mathrm{BD}}.
\end{equation}
Multiply by a test function \(w\) and integrate:
\begin{equation}
    \int_\Omega w\tau\phi_{tt}\,\dd V
    +\int_\Omega w\phi_t\,\dd V
    -\int_\Omega w\nabla\cdot(\alpha\nabla\phi)\,\dd V
    -\int_\Omega wS\,\dd V=0.
\end{equation}
Use Eq.~\eqref{eq:fem_parts_identity} on the diffusion term. The weak form becomes
\begin{equation}
    \boxed{
    \int_\Omega w\tau\phi_{tt}\,\dd V
    +\int_\Omega w\phi_t\,\dd V
    +\int_\Omega \nabla w\cdot\alpha\nabla\phi\,\dd V
    =\int_\Omega wS\,\dd V
    +\int_{\partial\Omega} w\alpha\nabla\phi\cdot\nvec\,\dd A.}
    \label{eq:fem_generic_weak}
\end{equation}
This is the central FEM statement for Module 1.

The term-by-term interpretation is:
\begin{align}
    \int_\Omega w\tau\phi_{tt}\,\dd V
    &:\quad \text{thermal inertia-like contribution from finite relaxation time},\\
    \int_\Omega w\phi_t\,\dd V
    &:\quad \text{relaxation/damping contribution},\\
    \int_\Omega \nabla w\cdot\alpha\nabla\phi\,\dd V
    &:\quad \text{diffusive coupling between neighboring regions},\\
    \int_\Omega wS\,\dd V
    &:\quad \text{volumetric source, including ballistic deposition},\\
    \int_{\partial\Omega} w\alpha\nabla\phi\cdot\nvec\,\dd A
    &:\quad \text{boundary flux contribution}.
\end{align}

\subsection{Conservative weak form for the ballistic source}
The source \(S_{\mathrm{BD}}\) contains \(-\nabla\cdot\qvec_b\). There are two ways to discretize it.

The direct way is
\begin{equation}
    F_{b,i}^{\mathrm{direct}}
    =\int_\Omega N_i\left(-\nabla\cdot\qvec_b\right)\,\dd V.
    \label{eq:direct_ballistic_load}
\end{equation}
This requires differentiating the ballistic flux field. If \(\qvec_b\) is computed by ray tracing or angular quadrature, numerical differentiation can be noisy.

A more conservative FEM form integrates the ballistic divergence by parts:
\begin{align}
    \int_\Omega w(-\nabla\cdot\qvec_b)\,\dd V
    &=-\int_{\partial\Omega} w\qvec_b\cdot\nvec\,\dd A
    +\int_\Omega \nabla w\cdot\qvec_b\,\dd V.
    \label{eq:ballistic_source_by_parts}
\end{align}
Thus the ballistic load can be assembled as
\begin{equation}
    \boxed{
    F_{b,i}
    =\int_\Omega \nabla N_i\cdot\qvec_b\,\dd V
    -\int_{\partial\Omega} N_i\qvec_b\cdot\nvec\,\dd A.}
    \label{eq:conservative_ballistic_load}
\end{equation}
This is usually better because it preserves the flux-balance interpretation: ballistic energy lost from one region appears as a source for the scattered medium component.

\section{Weak form for the one-dimensional validation equation}
For teaching and validation, let
\begin{equation}
    \phi(\eta,t^*)=\theta_m(\eta,t^*),
    \qquad
    D=\frac{\mathrm{Kn}^2}{3},
    \qquad
    s_b=-\mathrm{Kn}\frac{\partial q_b^*}{\partial\eta}.
\end{equation}
The strong form is
\begin{equation}
    \phi_{t^*t^*}+\phi_{t^*}-D\phi_{\eta\eta}=s_b.
\end{equation}
Multiply by \(w(\eta)\) and integrate from \(0\) to \(1\):
\begin{equation}
    \int_0^1 w\phi_{t^*t^*}\,\dd\eta
    +\int_0^1 w\phi_{t^*}\,\dd\eta
    -\int_0^1 wD\phi_{\eta\eta}\,\dd\eta
    =\int_0^1 ws_b\,\dd\eta.
\end{equation}
Integrate the second-derivative term by parts:
\begin{equation}
    -\int_0^1 wD\phi_{\eta\eta}\,\dd\eta
    =\int_0^1 D w_\eta\phi_\eta\,\dd\eta
    -\left[wD\phi_\eta\right]_{0}^{1}.
\end{equation}
Therefore
\begin{equation}
    \boxed{
    \int_0^1 w\phi_{t^*t^*}\,\dd\eta
    +\int_0^1 w\phi_{t^*}\,\dd\eta
    +\int_0^1 D w_\eta\phi_\eta\,\dd\eta
    =\int_0^1 ws_b\,\dd\eta
    +\left[wD\phi_\eta\right]_{0}^{1}.}
    \label{eq:fem_1d_weak}
\end{equation}
The bracketed boundary term means
\begin{equation}
    \left[wD\phi_\eta\right]_0^1
    =w(1)D\phi_\eta(1)-w(0)D\phi_\eta(0).
\end{equation}
This is the one-dimensional version of the general boundary integral \(\int_{\partial\Omega} w\alpha\nabla\phi\cdot\nvec\,\dd A\).

\section{The finite-dimensional approximation}
\subsection{Mesh and nodes}
Divide the one-dimensional domain \([0,1]\) into \(n_e\) finite elements. The nodes are
\begin{equation}
    0=\eta_1<\eta_2<\cdots<\eta_{n_n}=1,
\end{equation}
where \(n_n=n_e+1\) for a simple linear mesh. Element \(e\) connects two nodes:
\begin{equation}
    \Omega_e=[\eta_a,\eta_b],
    \qquad
    h_e=\eta_b-\eta_a.
\end{equation}
The FEM approximation is
\begin{equation}
    \boxed{
    \phi_h(\eta,t^*)=\sum_{j=1}^{n_n}\Phi_j(t^*)N_j(\eta).}
    \label{eq:fem_global_approx}
\end{equation}
The basis function \(N_j\) is usually chosen so that
\begin{equation}
    N_j(\eta_i)=\delta_{ij}.
\end{equation}
Therefore the coefficient \(\Phi_j\) is the approximate value of \(\phi\) at node \(j\):
\begin{equation}
    \Phi_j(t^*)=\phi_h(\eta_j,t^*).
\end{equation}

\subsection{Galerkin choice of test functions}
The Galerkin method uses the same functions for approximation and testing:
\begin{equation}
    w=N_i,
    \qquad i=1,2,\ldots,n_n.
\end{equation}
This gives one algebraic equation per node. The result is a matrix system for the nodal vector
\begin{equation}
    \bm{\Phi}(t^*)=
    \begin{bmatrix}
    \Phi_1(t^*) & \Phi_2(t^*) & \cdots & \Phi_{n_n}(t^*)
    \end{bmatrix}^{T}.
\end{equation}

\section{Linear shape functions in one dimension}
\subsection{Reference element}
A clean way to derive element matrices is to define a reference element \(\xi\in[-1,1]\). The two linear shape functions are
\begin{equation}
    \widehat N_1(\xi)=\frac{1-\xi}{2},
    \qquad
    \widehat N_2(\xi)=\frac{1+\xi}{2}.
    \label{eq:linear_shape_reference}
\end{equation}
They satisfy
\begin{equation}
    \widehat N_1(-1)=1,
    \quad
    \widehat N_1(1)=0,
    \qquad
    \widehat N_2(-1)=0,
    \quad
    \widehat N_2(1)=1.
\end{equation}
The physical coordinate inside element \(e\) is mapped by
\begin{equation}
    \eta(\xi)=\widehat N_1(\xi)\eta_a+\widehat N_2(\xi)\eta_b.
\end{equation}
The Jacobian of the map is
\begin{equation}
    J_e=\frac{\dd\eta}{\dd\xi}=\frac{h_e}{2}.
\end{equation}
Thus
\begin{equation}
    \dd\eta=J_e\,\dd\xi=\frac{h_e}{2}\,\dd\xi,
    \qquad
    \frac{\dd}{\dd\eta}=\frac{\dd\xi}{\dd\eta}\frac{\dd}{\dd\xi}=\frac{2}{h_e}\frac{\dd}{\dd\xi}.
\end{equation}

\subsection{Shape functions in physical coordinates}
On the physical element \([\eta_a,\eta_b]\), the same functions can be written as
\begin{equation}
    N_1^{(e)}(\eta)=\frac{\eta_b-\eta}{h_e},
    \qquad
    N_2^{(e)}(\eta)=\frac{\eta-\eta_a}{h_e}.
    \label{eq:linear_shape_physical}
\end{equation}
Their derivatives are constant:
\begin{equation}
    \frac{\dd N_1^{(e)}}{\dd\eta}=-\frac{1}{h_e},
    \qquad
    \frac{\dd N_2^{(e)}}{\dd\eta}=+\frac{1}{h_e}.
    \label{eq:linear_shape_derivatives}
\end{equation}
This is why linear one-dimensional FEM is so transparent: the field is piecewise linear and the gradient is piecewise constant.

\section{Element-level approximation}
On element \(e\), collect the two local nodal unknowns into
\begin{equation}
    \bm{\Phi}_e=
    \begin{bmatrix}
    \Phi_a\\
    \Phi_b
    \end{bmatrix}.
\end{equation}
The element interpolation is
\begin{equation}
    \phi_h^{(e)}(\eta,t^*)
    =N_1^{(e)}(\eta)\Phi_a(t^*)+N_2^{(e)}(\eta)\Phi_b(t^*)
    =\bm{N}_e(\eta)\bm{\Phi}_e(t^*),
\end{equation}
where
\begin{equation}
    \bm{N}_e=
    \begin{bmatrix}
    N_1^{(e)} & N_2^{(e)}
    \end{bmatrix}.
\end{equation}
The spatial derivative is
\begin{equation}
    \frac{\partial\phi_h^{(e)}}{\partial\eta}
    =\frac{\dd N_1^{(e)}}{\dd\eta}\Phi_a
    +\frac{\dd N_2^{(e)}}{\dd\eta}\Phi_b
    =\bm{B}_e\bm{\Phi}_e,
\end{equation}
with
\begin{equation}
    \bm{B}_e=
    \begin{bmatrix}
    -1/h_e & +1/h_e
    \end{bmatrix}.
\end{equation}

\section{Element mass matrix}
The first two weak-form terms involve \(\phi_{t^*t^*}\) and \(\phi_{t^*}\). Since the shape functions are time-independent,
\begin{equation}
    \phi_{h,t^*}^{(e)}=\bm{N}_e\dot{\bm{\Phi}}_e,
    \qquad
    \phi_{h,t^*t^*}^{(e)}=\bm{N}_e\ddot{\bm{\Phi}}_e.
\end{equation}
Choosing \(w=N_i\), the element contribution is
\begin{equation}
    \int_{\Omega_e} N_i\phi_{h,t^*t^*}^{(e)}\,\dd\eta
    =\sum_{j=1}^{2}\left(\int_{\Omega_e}N_iN_j\,\dd\eta\right)\ddot\Phi_j.
\end{equation}
This defines the element mass matrix
\begin{equation}
    \boxed{M_{e,ij}=\int_{\Omega_e}N_iN_j\,\dd\eta.}
    \label{eq:element_mass_def}
\end{equation}
For a linear element, compute the entries explicitly. With \(x\) temporarily denoting distance from the left element node, \(0\le x\le h_e\),
\begin{equation}
    N_1=1-\frac{x}{h_e},
    \qquad
    N_2=\frac{x}{h_e}.
\end{equation}
Then
\begin{align}
    M_{e,11}&=\int_0^{h_e}\left(1-\frac{x}{h_e}\right)^2\dd x=\frac{h_e}{3},\\
    M_{e,22}&=\int_0^{h_e}\left(\frac{x}{h_e}\right)^2\dd x=\frac{h_e}{3},\\
    M_{e,12}&=M_{e,21}=\int_0^{h_e}\left(1-\frac{x}{h_e}\right)\left(\frac{x}{h_e}\right)\dd x=\frac{h_e}{6}.
\end{align}
Therefore
\begin{equation}
    \boxed{
    M_e=\frac{h_e}{6}
    \begin{bmatrix}
    2 & 1\\
    1 & 2
    \end{bmatrix}.}
    \label{eq:linear_element_mass}
\end{equation}
This matrix appears twice in the ballistic-diffusive system: once multiplying \(\ddot{\bm{\Phi}}\) and once multiplying \(\dot{\bm{\Phi}}\). In dimensional form the second-derivative term is multiplied by \(\tau\), so the inertial mass-like matrix is \(\tau M\).

\section{Element stiffness matrix}
The diffusion term in the weak form is
\begin{equation}
    \int_{\Omega_e}D\frac{\dd N_i}{\dd\eta}\frac{\dd\phi_h}{\dd\eta}\,\dd\eta
    =\sum_{j=1}^{2}\left(\int_{\Omega_e}D\frac{\dd N_i}{\dd\eta}\frac{\dd N_j}{\dd\eta}\,\dd\eta\right)\Phi_j.
\end{equation}
This defines
\begin{equation}
    \boxed{K_{e,ij}=\int_{\Omega_e}D\frac{\dd N_i}{\dd\eta}\frac{\dd N_j}{\dd\eta}\,\dd\eta.}
    \label{eq:element_stiffness_def}
\end{equation}
Using Eq.~\eqref{eq:linear_shape_derivatives},
\begin{align}
    K_{e,11}&=\int_{\Omega_e}D\left(-\frac{1}{h_e}\right)\left(-\frac{1}{h_e}\right)\dd\eta=\frac{D}{h_e},\\
    K_{e,12}&=\int_{\Omega_e}D\left(-\frac{1}{h_e}\right)\left(+\frac{1}{h_e}\right)\dd\eta=-\frac{D}{h_e},\\
    K_{e,21}&=-\frac{D}{h_e},\\
    K_{e,22}&=\frac{D}{h_e}.
\end{align}
So
\begin{equation}
    \boxed{
    K_e=\frac{D}{h_e}
    \begin{bmatrix}
    1 & -1\\
    -1 & 1
    \end{bmatrix},
    \qquad
    D=\frac{\mathrm{Kn}^2}{3}.}
    \label{eq:linear_element_stiffness}
\end{equation}

\begin{physicsbox}
\textbf{Physical meaning of stiffness.} The vector \(K_e\bm{\Phi}_e\) penalizes differences between neighboring nodal values. If \(\Phi_a=\Phi_b\), the element has no gradient and \(K_e\bm{\Phi}_e=\bm{0}\). Diffusion only acts when there is spatial variation.
\end{physicsbox}

\section{Element source vector}
\subsection{Direct ballistic source vector}
The source vector for a known scalar source \(s_b(\eta,t^*)\) is
\begin{equation}
    \boxed{F_{e,i}^{(s)}(t^*)=\int_{\Omega_e}N_i(\eta)s_b(\eta,t^*)\,\dd\eta.}
    \label{eq:element_source_def}
\end{equation}
If \(s_b\) is approximately constant inside the element, \(s_b\approx s_e\), then
\begin{equation}
    F_e^{(s)}\approx s_e\int_{\Omega_e}
    \begin{bmatrix}
    N_1\\N_2
    \end{bmatrix}\dd\eta
    =s_e\frac{h_e}{2}
    \begin{bmatrix}
    1\\1
    \end{bmatrix}.
    \label{eq:constant_source_vector}
\end{equation}
This is the simplest load-vector formula.

\subsection{Conservative ballistic source vector in one dimension}
Because
\begin{equation}
    s_b=-\mathrm{Kn}\frac{\partial q_b^*}{\partial\eta},
\end{equation}
we can write
\begin{align}
    F_{e,i}^{(b)}
    &=-\mathrm{Kn}\int_{\Omega_e}N_i\frac{\partial q_b^*}{\partial\eta}\,\dd\eta\\
    &=-\mathrm{Kn}\left[N_iq_b^*\right]_{\eta_a}^{\eta_b}
    +\mathrm{Kn}\int_{\Omega_e}\frac{\dd N_i}{\dd\eta}q_b^*\,\dd\eta.
    \label{eq:conservative_1d_ballistic_load}
\end{align}
This form uses \(q_b^*\) itself rather than a numerical derivative of \(q_b^*\). For ray-based or quadrature-based ballistic flux, this is usually more stable and more faithful to energy conservation.

For a linear element, the derivatives \(\dd N_i/\dd\eta\) are constants, so the volume part is especially simple:
\begin{equation}
    \mathrm{Kn}\int_{\Omega_e}\frac{\dd N_1}{\dd\eta}q_b^*\,\dd\eta
    =-\frac{\mathrm{Kn}}{h_e}\int_{\Omega_e}q_b^*\,\dd\eta,
\end{equation}
\begin{equation}
    \mathrm{Kn}\int_{\Omega_e}\frac{\dd N_2}{\dd\eta}q_b^*\,\dd\eta
    =+\frac{\mathrm{Kn}}{h_e}\int_{\Omega_e}q_b^*\,\dd\eta.
\end{equation}
The boundary part transfers ballistic flux across element faces. When assembled over all elements, interior face contributions cancel in pairs, leaving only physical domain-boundary ballistic flux plus elementwise deposition due to attenuation.

\section{Numerical quadrature}
For simple linear elements with constant \(D\), the mass and stiffness matrices can be integrated analytically. In general, FEM codes use quadrature. On the reference element,
\begin{equation}
    \int_{\Omega_e}g(\eta)\,\dd\eta
    =\int_{-1}^{1}g[\eta(\xi)]J_e\,\dd\xi
    \approx\sum_{q=1}^{n_q}w_q g[\eta(\xi_q)]J_e.
\end{equation}
For two-point Gauss quadrature on \([-1,1]\),
\begin{equation}
    \xi_1=-\frac{1}{\sqrt{3}},
    \qquad
    \xi_2=+\frac{1}{\sqrt{3}},
    \qquad
    w_1=w_2=1.
\end{equation}
This quadrature exactly integrates polynomials up to degree three. Since the linear-element mass integrand \(N_iN_j\) is quadratic, two-point Gauss quadrature exactly recovers Eq.~\eqref{eq:linear_element_mass}.

For the ballistic source, quadrature is not merely a convenience. The function \(q_b^*(\eta,t^*)\) can contain sharp arrival fronts because of the finite-arrival condition \(\mu_t=\eta/(\mathrm{Kn}t^*)\). A practical implementation should evaluate \(q_b^*\) at quadrature points and use enough spatial and temporal resolution to capture those fronts.

\section{Global assembly}
Each element produces a local matrix equation of the form
\begin{equation}
    M_e\ddot{\bm{\Phi}}_e+M_e\dot{\bm{\Phi}}_e+K_e\bm{\Phi}_e=\bm{F}_e.
\end{equation}
The local vector \(\bm{\Phi}_e\) contains only the degrees of freedom attached to the element. The global vector \(\bm{\Phi}\) contains all nodal degrees of freedom. Assembly is the process of adding local contributions into the correct global rows and columns.

Suppose element \(e\) connects global nodes \(a\) and \(b\). Its local-to-global map is
\begin{equation}
    \mathrm{conn}(e,:)=[a,b].
\end{equation}
Then the assembly rule is
\begin{equation}
    M_{A B}\leftarrow M_{A B}+M_{e,ij},
    \qquad
    K_{A B}\leftarrow K_{A B}+K_{e,ij},
    \qquad
    F_A\leftarrow F_A+F_{e,i},
\end{equation}
where
\begin{equation}
    A=\mathrm{conn}(e,i),
    \qquad
    B=\mathrm{conn}(e,j).
\end{equation}

For three linear elements, the stiffness matrix has the schematic tridiagonal form
\begin{equation}
    K=\begin{bmatrix}
    k_1 & -k_1 & 0 & 0\\
    -k_1 & k_1+k_2 & -k_2 & 0\\
    0 & -k_2 & k_2+k_3 & -k_3\\
    0 & 0 & -k_3 & k_3
    \end{bmatrix},
    \qquad
    k_e=\frac{D}{h_e}.
\end{equation}
The diagonal entries collect all elements touching a node. The off-diagonal entries couple neighboring nodes.

\begin{notebox}
\textbf{Implementation meaning.} Assembly is not mysterious. It is just bookkeeping: compute each small element matrix, then add its numbers into the large global sparse matrix at the rows and columns corresponding to that element's nodes.
\end{notebox}

\section{Boundary conditions in the FEM system}
\subsection{Natural Neumann boundary condition}
From the weak form, the natural boundary term is
\begin{equation}
    \int_{\partial\Omega}w\alpha\nabla\phi\cdot\nvec\,\dd A.
\end{equation}
Define the outward diffusive energy flux as
\begin{equation}
    g=-\alpha\nabla\phi\cdot\nvec.
\end{equation}
Then
\begin{equation}
    \alpha\nabla\phi\cdot\nvec=-g,
\end{equation}
and the boundary contribution is
\begin{equation}
    -\int_{\Gamma_N}w g\,\dd A.
\end{equation}
Thus, depending on sign convention, a prescribed outward heat flux may enter the load vector with a minus sign. The safest rule is to start from the weak form and explicitly define whether the prescribed flux is \(\alpha\nabla\phi\cdot\nvec\) or \(-\alpha\nabla\phi\cdot\nvec\).

In one dimension, the boundary contribution is
\begin{equation}
    \left[wD\phi_\eta\right]_0^1=w(1)D\phi_\eta(1)-w(0)D\phi_\eta(0).
\end{equation}
At \(\eta=1\), the outward normal is \(+1\). At \(\eta=0\), the outward normal is \(-1\). That is why the two endpoint signs differ.

\subsection{Dirichlet boundary condition}
A Dirichlet boundary condition prescribes
\begin{equation}
    \phi=\bar\phi\quad \text{on }\Gamma_D.
\end{equation}
If node \(p\) is Dirichlet, then
\begin{equation}
    \Phi_p=\bar\phi_p.
\end{equation}
The clean elimination method partitions the unknowns into free and prescribed components:
\begin{equation}
    \bm{\Phi}=\begin{bmatrix}\bm{\Phi}_f\\\bm{\Phi}_d\end{bmatrix}.
\end{equation}
For a static-looking linear solve \(A\bm{\Phi}=\bm{b}\), the free equations become
\begin{equation}
    A_{ff}\bm{\Phi}_f=\bm{b}_f-A_{fd}\bm{\Phi}_d.
\end{equation}
This method preserves symmetry of \(A_{ff}\). Another common coding method is row-column modification: zero the constrained row and column, put a one on the diagonal, and set the right-hand side to the prescribed value.

In the ballistic-diffusive model, one must be careful with Dirichlet temperature boundaries. If the wall emission has already been represented by \(\qvec_b\), then prescribing the same wall temperature directly on \(\theta_m\) can double-count boundary physics. For validation, Dirichlet conditions are fine; for the physical ballistic-diffusive model, Marshak-type boundaries are more consistent.

\subsection{Marshak/Robin boundary condition}
The medium boundary condition derived earlier is
\begin{equation}
    \tau\phi_t+\phi=-\ell_M\nabla\phi\cdot\nvec,
    \qquad
    \ell_M=\frac{\mfp}{3}
    \label{eq:marshak_general_phi}
\end{equation}
in dimensional notation. Rearranging gives
\begin{equation}
    \nabla\phi\cdot\nvec=-\frac{1}{\ell_M}(\tau\phi_t+\phi).
\end{equation}
Therefore the weak-form boundary term is
\begin{equation}
    \int_{\Gamma_R}w\alpha\nabla\phi\cdot\nvec\,\dd A
    =-\int_{\Gamma_R}w\frac{\alpha}{\ell_M}(\tau\phi_t+\phi)\,\dd A.
\end{equation}
Move this contribution to the left side:
\begin{equation}
    \int_{\Omega} w\tau\phi_{tt}\,\dd V
    +\int_{\Omega} w\phi_t\,\dd V
    +\int_{\Omega}\nabla w\cdot\alpha\nabla\phi\,\dd V
    +\int_{\Gamma_R}w\frac{\alpha\tau}{\ell_M}\phi_t\,\dd A
    +\int_{\Gamma_R}w\frac{\alpha}{\ell_M}\phi\,\dd A
    =\int_\Omega wS\,\dd V.
\end{equation}
This shows that the Marshak boundary adds two boundary matrices:
\begin{align}
    C^{(R)}_{ij} &=\int_{\Gamma_R}N_i\frac{\alpha\tau}{\ell_M}N_j\,\dd A,\\
    K^{(R)}_{ij} &=\int_{\Gamma_R}N_i\frac{\alpha}{\ell_M}N_j\,\dd A.
\end{align}
The first modifies the damping matrix, and the second modifies the stiffness matrix.

For the nondimensional one-dimensional equation,
\begin{equation}
    \phi_{t^*}+\phi=-\frac{\mathrm{Kn}}{3}\frac{\partial\phi}{\partial n},
\end{equation}
and \(D=\mathrm{Kn}^2/3\). Thus
\begin{equation}
    D\frac{\partial\phi}{\partial n}
    =-\mathrm{Kn}(\phi_{t^*}+\phi).
\end{equation}
The boundary term becomes
\begin{equation}
    \int_{\partial\Omega}wD\frac{\partial\phi}{\partial n}\,\dd\Gamma
    =-\int_{\partial\Omega}w\mathrm{Kn}(\phi_{t^*}+\phi)\,\dd\Gamma.
\end{equation}
Moving it to the left adds
\begin{equation}
    \mathrm{Kn}\int_{\partial\Omega}w\phi_{t^*}\,\dd\Gamma
    +\mathrm{Kn}\int_{\partial\Omega}w\phi\,\dd\Gamma.
\end{equation}
In a one-dimensional endpoint discretization, this produces contributions at the endpoint nodes. For example, if both endpoints are Marshak boundaries, the boundary matrix is schematically
\begin{equation}
    H=\mathrm{Kn}
    \begin{bmatrix}
    1 & 0 & \cdots & 0\\
    0 & 0 & \cdots & 0\\
    \vdots & \vdots & \ddots & \vdots\\
    0 & 0 & \cdots & 1
    \end{bmatrix}.
\end{equation}
Then the semidiscrete equation becomes
\begin{equation}
    \boxed{
    M\ddot{\bm{\Phi}}+(M+H)\dot{\bm{\Phi}}+(K+H)\bm{\Phi}=\bm{F}_b.}
    \label{eq:marshak_semidiscrete_1d}
\end{equation}
This equation is often the most direct form for the 1D validation code.

\section{Global semidiscrete matrix equation}
Ignoring boundary matrices for the moment, the global FEM system is
\begin{equation}
    \boxed{
    M\ddot{\bm{\Phi}}+M\dot{\bm{\Phi}}+K\bm{\Phi}=\bm{F}_b(t^*)+\bm{F}_{\mathrm{ext}}(t^*)+\bm{F}_{\partial\Omega}(t^*).}
    \label{eq:fem_global_1d_matrix}
\end{equation}
For the dimensional equation, this becomes
\begin{equation}
    \boxed{
    \tau M\ddot{\bm{U}}+M\dot{\bm{U}}+K\bm{U}=\bm{F}_{\mathrm{BD}}+\bm{F}_{\partial\Omega}.}
    \label{eq:fem_dimensional_matrix_full}
\end{equation}
Here
\begin{align}
    M_{ij} &=\int_\Omega N_iN_j\,\dd V,\\
    K_{ij} &=\int_\Omega \nabla N_i\cdot\alpha\nabla N_j\,\dd V,\\
    (F_{\mathrm{BD}})_i &=\int_\Omega N_iS_{\mathrm{BD}}\,\dd V,\\
    (F_{\partial\Omega})_i &=\int_{\partial\Omega}N_i\alpha\nabla\phi\cdot\nvec\,\dd A.
\end{align}
With Marshak boundaries moved to the left, the form is better written as
\begin{equation}
    \boxed{
    \tau M\ddot{\bm{U}}+(M+C_R)\dot{\bm{U}}+(K+K_R)\bm{U}=\bm{F}_{\mathrm{BD}}.}
    \label{eq:fem_dimensional_matrix_robin}
\end{equation}
This is the FEM analog of a damped second-order system. The analogy is useful but should not be overinterpreted: \(\bm{U}\) is not a mechanical displacement. It is nodal medium energy.

\section{Time discretization and solving the nodal values}
After spatial FEM, the problem is still continuous in time. We now need nodal values at discrete time levels
\begin{equation}
    t_n=n\Delta t,
    \qquad
    \bm{\Phi}^n\approx \bm{\Phi}(t_n).
\end{equation}
For clarity, consider the nondimensional semidiscrete equation
\begin{equation}
    M\ddot{\bm{\Phi}}+C\dot{\bm{\Phi}}+K\bm{\Phi}=\bm{F}^{n},
    \label{eq:fem_second_order_generic}
\end{equation}
where \(C=M\) without Marshak boundaries and \(C=M+H\) with nondimensional Marshak boundaries.

\subsection{Backward-difference update}
A simple implicit second-order finite-difference approximation is
\begin{equation}
    \ddot{\bm{\Phi}}^{n+1}\approx
    \frac{\bm{\Phi}^{n+1}-2\bm{\Phi}^{n}+\bm{\Phi}^{n-1}}{\Delta t^2},
    \qquad
    \dot{\bm{\Phi}}^{n+1}\approx
    \frac{\bm{\Phi}^{n+1}-\bm{\Phi}^{n}}{\Delta t}.
\end{equation}
Substitute into Eq.~\eqref{eq:fem_second_order_generic} at \(t_{n+1}\):
\begin{equation}
    M\frac{\bm{\Phi}^{n+1}-2\bm{\Phi}^{n}+\bm{\Phi}^{n-1}}{\Delta t^2}
    +C\frac{\bm{\Phi}^{n+1}-\bm{\Phi}^{n}}{\Delta t}
    +K\bm{\Phi}^{n+1}
    =\bm{F}^{n+1}.
\end{equation}
Collect unknown \(\bm{\Phi}^{n+1}\) terms on the left:
\begin{equation}
    \boxed{
    A_{\mathrm{eff}}\bm{\Phi}^{n+1}=\bm{b}^{n+1},}
    \label{eq:effective_solve}
\end{equation}
where
\begin{equation}
    \boxed{
    A_{\mathrm{eff}}=\frac{1}{\Delta t^2}M+\frac{1}{\Delta t}C+K,}
    \label{eq:aeff_bdf}
\end{equation}
and
\begin{equation}
    \boxed{
    \bm{b}^{n+1}=\bm{F}^{n+1}
    +M\left(\frac{2\bm{\Phi}^{n}-\bm{\Phi}^{n-1}}{\Delta t^2}\right)
    +C\left(\frac{\bm{\Phi}^{n}}{\Delta t}\right).}
    \label{eq:rhs_bdf}
\end{equation}
This is the concrete meaning of ``solving the nodal values.'' At every time step, assemble or update the right-hand side and solve a sparse linear system for the next nodal vector.

For the dimensional system \(\tau M\ddot{\bm{U}}+C\dot{\bm{U}}+K\bm{U}=\bm{F}\), replace \(M\) in the acceleration contribution by \(\tau M\):
\begin{equation}
    A_{\mathrm{eff}}=\frac{\tau}{\Delta t^2}M+\frac{1}{\Delta t}C+K.
\end{equation}

\subsection{Initial conditions}
The second-order-in-time equation requires two initial conditions:
\begin{equation}
    \bm{\Phi}^{0}=\bm{\Phi}(0),
    \qquad
    \dot{\bm{\Phi}}^{0}=\dot{\bm{\Phi}}(0).
\end{equation}
For the canonical thin-film validation,
\begin{equation}
    \bm{\Phi}^{0}=\bm{0},
    \qquad
    \dot{\bm{\Phi}}^{0}=\bm{0}
\end{equation}
for the perturbation medium energy. A backward-difference scheme also needs \(\bm{\Phi}^{-1}\). A first-order Taylor approximation gives
\begin{equation}
    \bm{\Phi}^{-1}=\bm{\Phi}^{0}-\Delta t\dot{\bm{\Phi}}^{0}.
\end{equation}
If both are zero, \(\bm{\Phi}^{-1}=\bm{0}\). For higher accuracy, the initial acceleration can be obtained from the semidiscrete equation:
\begin{equation}
    M\ddot{\bm{\Phi}}^{0}=\bm{F}^{0}-C\dot{\bm{\Phi}}^{0}-K\bm{\Phi}^{0},
\end{equation}
and then
\begin{equation}
    \bm{\Phi}^{-1}=\bm{\Phi}^{0}-\Delta t\dot{\bm{\Phi}}^{0}+\frac{\Delta t^2}{2}\ddot{\bm{\Phi}}^{0}.
\end{equation}

\subsection{Algorithmic summary}
A minimal 1D FEM solver follows this sequence:
\begin{enumerate}[leftmargin=1.6em]
    \item Create the nodal mesh \(\eta_i\) and element connectivity.
    \item For each element, compute \(M_e\) and \(K_e\).
    \item Assemble global sparse matrices \(M\) and \(K\).
    \item Add boundary matrices \(H\) if using Marshak/Robin boundaries.
    \item Initialize \(\bm{\Phi}^0\), \(\dot{\bm{\Phi}}^0\), and \(\bm{\Phi}^{-1}\).
    \item At every time step, compute the ballistic flux \(q_b^*(\eta,t^*)\) and load vector \(\bm{F}_b^{n+1}\).
    \item Form \(A_{\mathrm{eff}}\) and \(\bm{b}^{n+1}\).
    \item Apply any essential Dirichlet constraints, if used.
    \item Solve \(A_{\mathrm{eff}}\bm{\Phi}^{n+1}=\bm{b}^{n+1}\).
    \item Reconstruct \(\theta_m^h(\eta,t^*)\), compute \(\theta_b\), and form \(\theta=\theta_m+\theta_b\).
\end{enumerate}

\section{Field reconstruction after solving nodes}
Once the nodal vector has been found, the FEM solution is not only a list of numbers. It defines a continuous piecewise-polynomial field:
\begin{equation}
    \boxed{
    \phi_h(\eta,t_n)=\sum_{j=1}^{n_n}\Phi_j^nN_j(\eta).}
\end{equation}
Inside element \(e=[\eta_a,\eta_b]\), this is simply
\begin{equation}
    \phi_h^{(e)}(\eta,t_n)=
    \frac{\eta_b-\eta}{h_e}\Phi_a^n+
    \frac{\eta-\eta_a}{h_e}\Phi_b^n.
\end{equation}
The gradient is
\begin{equation}
    \frac{\partial\phi_h^{(e)}}{\partial\eta}
    =\frac{\Phi_b^n-\Phi_a^n}{h_e}.
\end{equation}
The scattered-medium diffusive flux in dimensionless variables is proportional to
\begin{equation}
    q_m^*=-D\frac{\partial\theta_m}{\partial\eta},
\end{equation}
while the total equivalent response plotted for validation is
\begin{equation}
    \boxed{
    \theta_h(\eta,t^*)=\theta_{m,h}(\eta,t^*)+\theta_b(\eta,t^*).}
\end{equation}
Thus the final approximation is not just the medium FEM field. The ballistic-diffusive model reconstructs the observable energy-temperature response by adding the explicitly computed ballistic energy contribution.

\section{Extension from 1D elements to 2D and 3D elements}
The same FEM logic generalizes to realistic chip geometry. For a multidimensional element, write
\begin{equation}
    \phi_h^{(e)}(\rvec,t)=\sum_{a=1}^{n_{en}}N_a^{(e)}(\rvec)\Phi_a(t),
\end{equation}
where \(n_{en}\) is the number of nodes per element. In matrix notation,
\begin{equation}
    \phi_h^{(e)}=\bm{N}_e\bm{\Phi}_e,
    \qquad
    \nabla\phi_h^{(e)}=\bm{B}_e\bm{\Phi}_e.
\end{equation}
The element matrices are
\begin{align}
    M_e &=\int_{\Omega_e}\bm{N}_e^{T}\bm{N}_e\,\dd V,\\
    K_e &=\int_{\Omega_e}\bm{B}_e^{T}\alpha\bm{B}_e\,\dd V,\\
    \bm{F}_e &=\int_{\Omega_e}\bm{N}_e^{T}S\,\dd V
    +\int_{\partial\Omega_e\cap\partial\Omega}\bm{N}_e^{T}\alpha\nabla\phi\cdot\nvec\,\dd A.
\end{align}
For an anisotropic material, \(\alpha\) becomes a tensor:
\begin{equation}
    \bm{\alpha}=\frac{\bm{k}}{C},
\end{equation}
and the stiffness matrix becomes
\begin{equation}
    K_e=\int_{\Omega_e}\bm{B}_e^{T}\bm{\alpha}\bm{B}_e\,\dd V.
\end{equation}
This matters for crystalline substrates and layered device stacks where heat transport can differ in-plane and through-thickness.

\subsection{Isoparametric mapping and Jacobian}
For quadrilateral, hexahedral, triangular, or tetrahedral elements, integrals are usually evaluated on a reference element with natural coordinates. In 3D, for example,
\begin{equation}
    \rvec(\xi,\zeta,\chi)=\sum_{a=1}^{n_{en}}N_a(\xi,\zeta,\chi)\rvec_a.
\end{equation}
The Jacobian matrix is
\begin{equation}
    \bm{J}=\frac{\partial(x,y,z)}{\partial(\xi,\zeta,\chi)}.
\end{equation}
Derivatives transform by
\begin{equation}
    \nabla_{\rvec}N_a=\bm{J}^{-T}\nabla_{\xi}N_a.
\end{equation}
The volume element transforms as
\begin{equation}
    \dd V=\abs{\det\bm{J}}\,\dd\xi\,\dd\zeta\,\dd\chi.
\end{equation}
Therefore the stiffness matrix is computed numerically as
\begin{equation}
    K_e\approx\sum_{q=1}^{n_q}\bm{B}_e(\xi_q)^T\bm{\alpha}(\xi_q)\bm{B}_e(\xi_q)
    \abs{\det\bm{J}(\xi_q)}w_q.
\end{equation}
This is the same formula as in 1D, only with more geometry bookkeeping.

\section{FEM-specific interpretation for ballistic-diffusive conduction}
The final FEM system tells a clear physical story:
\begin{equation}
\begin{gathered}
    \text{thermal inertia}+\text{relaxation}+\text{diffusive redistribution}\\
    =\text{ballistic deposition}+\text{external heating}+\text{boundary exchange}.
\end{gathered}
\end{equation}
In matrix form,
\begin{equation}
    \tau M\ddot{\bm{U}}+(M+C_R)\dot{\bm{U}}+(K+K_R)\bm{U}=\bm{F}_{b}+\bm{F}_{e}.
\end{equation}
The terms should be read as follows:
\begin{align}
    \tau M\ddot{\bm{U}} &: \text{finite relaxation time; prevents purely parabolic instantaneous response},\\
    M\dot{\bm{U}} &: \text{relaxation of medium energy},\\
    K\bm{U} &: \text{spatial smoothing by scattered-carrier diffusion},\\
    C_R\dot{\bm{U}}+K_R\bm{U} &: \text{Marshak boundary exchange of the medium component},\\
    \bm{F}_b &: \text{energy transferred from ballistic boundary-originating carriers into the medium},\\
    \bm{F}_e &: \text{external volumetric heating and its relaxation correction}.
\end{align}

\begin{physicsbox}
\textbf{Most important implementation warning.} Do not treat the ballistic-diffusive FEM solve as an ordinary heat equation with a different conductivity. The unknown solved by FEM is the medium/scattered energy. The ballistic part is computed separately and enters both as a source term and as an additive contribution to the reconstructed total energy response.
\end{physicsbox}

\section{Common checks for the Module 1 FEM implementation}
A correct FEM implementation should satisfy the following checks.

\subsection{Patch test for constant field}
If \(\phi\) is constant and all sources are zero, then
\begin{equation}
    K\bm{\Phi}=\bm{0}.
\end{equation}
This follows from the fact that constant fields have zero gradient. Failure of this test usually indicates an error in shape-function derivatives, Jacobian factors, or assembly.

\subsection{Symmetry and positivity}
For positive \(D\) or positive \(\alpha\), the stiffness matrix should be symmetric positive semidefinite before Dirichlet constraints:
\begin{equation}
    \bm{x}^TK\bm{x}\ge 0.
\end{equation}
The mass matrix should be symmetric positive definite for unconstrained physical nodes:
\begin{equation}
    \bm{x}^TM\bm{x}>0
    \quad \text{for }\bm{x}\ne\bm{0}.
\end{equation}
These are structural checks on the FEM discretization.

\subsection{Energy-balance check}
For the conservative ballistic loading, integrating over the whole domain should reproduce the net ballistic flux entering minus leaving, plus attenuation/deposition terms. Numerically, this means that the sum of nodal source-vector entries should agree with the domain-integrated source:
\begin{equation}
    \sum_i F_{b,i}\approx\int_\Omega -\nabla\cdot\qvec_b\,\dd V.
\end{equation}
If the conservative form Eq.~\eqref{eq:conservative_ballistic_load} is implemented correctly, this check is much more robust than differentiating a noisy \(q_b\) field.

\subsection{Diffusive-limit check}
As \(\mathrm{Kn}\ll1\), after the earliest transient, the solution should approach Fourier-like behavior. This does not mean the ballistic-diffusive equation literally becomes the same matrix equation at all times. It means the computed total response should lose strong boundary-memory effects when the mean free path is much smaller than the film thickness.

\subsection{Arrival-time check}
For \(\mathrm{Kn}\sim1\) or larger, the ballistic contribution must show delayed arrival. At position \(\eta\), no left-boundary ballistic contribution should appear before
\begin{equation}
    t^*<\frac{\eta}{\mathrm{Kn}}
\end{equation}
for the fastest forward carriers with \(\mu=1\). If the code shows immediate ballistic response everywhere, the finite-arrival lower limit in the angular integral has been implemented incorrectly.

\section{Practical implementation strategy for this project}
The project should not begin with the full qubit geometry. The physically defensible sequence is:
\begin{enumerate}[leftmargin=1.6em]
    \item Validate the one-dimensional thin-film equation for several \(\mathrm{Kn}\) values.
    \item Confirm the diffusive limit as \(\mathrm{Kn}\ll1\).
    \item Confirm finite propagation and reduced early-time heat flux as \(\mathrm{Kn}\sim1\).
    \item Add a simple three-dimensional box before using a realistic transmon geometry.
    \item Only then couple the transport model to geometry-dependent trap and heat-sink studies.
\end{enumerate}

For Module 1, the primary code target is therefore
\begin{equation}
    \texttt{main\_module1\_bd\_1d\_validation.m}.
\end{equation}
It should generate temperature/internal-energy profiles, ballistic flux terms, and surface flux histories for
\begin{equation}
    \mathrm{Kn}=0.01,\ 0.1,\ 1,\ 10.
\end{equation}

\section{How this connects to the superconducting-qubit project}
The final goal of this project is not simply to solve a heat equation. The larger goal is to understand how geometry and materials influence energy flow and quasiparticle poisoning in superconducting qubit structures. Module 1 supplies the thermal transport language needed for later modules.

In a qubit device, local energy can be deposited near a Josephson junction, in a capacitor pad, in a substrate, or at an interface. The early-time flow of that energy can determine whether high-energy phonons or quasiparticles reach the sensitive junction region. A trap or heat sink changes the geometry and boundary conditions of the transport problem. Therefore, a model that captures finite carrier flight, boundary memory, and diffusion after scattering is more appropriate than immediately assuming ordinary Fourier conduction everywhere.

\section{Assumptions and limitations}
The ballistic-diffusive model used here is still an approximation. Its main assumptions are:
\begin{enumerate}[leftmargin=1.6em]
    \item \textbf{Particle-like carriers.} The model assumes a particle transport picture. It does not describe coherent wave interference of phonons.
    \item \textbf{Relaxation-time approximation.} Detailed scattering physics is compressed into \(\tau_\omega\) or an averaged \(\tau\).
    \item \textbf{Angular P1 closure for scattered carriers.} The medium component is approximated as nearly isotropic plus first-order anisotropy.
    \item \textbf{Averaged material properties.} The first implementation uses effective \(C\), \(k\), \(v\), \(\tau\), and \(\mfp\), rather than fully mode-resolved phonon bands.
    \item \textbf{Boundary modeling matters.} The ballistic source depends strongly on emitted/reflected boundary intensity. A poor boundary model can dominate the result.
    \item \textbf{Room-temperature foundation.} Cryogenic superconducting changes, quasiparticle generation, recombination, and electron-phonon nonequilibrium are deferred to later modules.
\end{enumerate}

\section{Summary of the derivation}
The derivation can be summarized as follows:
\begin{enumerate}[leftmargin=1.6em]
    \item Start from local energy conservation: \(\partial u/\partial t+\nabla\cdot\qvec=\dot{q}_e\).
    \item Describe heat carriers using the BTE under RTA: \(\partial f/\partial t+\bm{v}\cdot\nabla f=-(f-f_0)/\tau\).
    \item Convert to intensity notation so that angular moments produce \(u\) and \(\qvec\).
    \item Split the intensity into \(I_\omega=I_{b\omega}+I_{m\omega}\).
    \item Treat \(I_{b\omega}\) by characteristics, giving retarded boundary emission and exponential attenuation.
    \item Treat \(I_{m\omega}\) using the P1 diffusion approximation.
    \item Derive a Cattaneo-like relation for only the medium flux: \(\tau\partial\qvec_m/\partial t+\qvec_m=-(k/C)\nabla u_m\).
    \item Combine with energy conservation to obtain
    \begin{equation*}
    \tau\frac{\partial^2u_m}{\partial t^2}+\frac{\partial u_m}{\partial t}
    =\nabla\cdot\left(\frac{k}{C}\nabla u_m\right)-\nabla\cdot\qvec_b
    +\left(\dot{q}_e+\tau\frac{\partial\dot{q}_e}{\partial t}\right).
    \end{equation*}
    \item Convert this equation into a FEM weak form for numerical implementation.
\end{enumerate}


\section{Lecture-ready slide translation for Module 1}
This section converts the derivation into presentation logic. Each item corresponds to one possible slide or slide cluster.

\begin{notebox}
\textbf{Slide 1: Why ordinary heat conduction breaks down.}
Physical question: When does \(\qvec=-k\nabla T\) stop being enough? Governing idea: Fourier conduction assumes local equilibrium and local flux response. Approximation being challenged: heat flux is determined only by the local temperature gradient. Figure to show: a chip-scale length \(L\) compared with carrier mean free path \(\mfp\).
\end{notebox}

\begin{notebox}
\textbf{Slide 2: Heat carriers and mean free path.}
Physical question: What actually carries thermal energy? Governing idea: phonons or other carriers move with group velocity \(v\) and scatter after time \(\tau\). Key equation: \(\mfp=v\tau\). Figure to show: carrier trajectories with short and long mean free paths.
\end{notebox}

\begin{notebox}
\textbf{Slide 3: The distribution function.}
Physical question: What replaces temperature when the system is not locally equilibrated? Governing idea: use \(f(t,\rvec,\bm{k})\), a phase-space occupation. Approximation: local equilibrium is represented by \(f_0\), but the actual distribution may be far from \(f_0\). Figure to show: isotropic equilibrium distribution versus directional nonequilibrium distribution.
\end{notebox}

\begin{notebox}
\textbf{Slide 4: Boltzmann transport equation.}
Physical question: How does the microscopic carrier distribution evolve? Governing equation: \(\partial f/\partial t+\bm{v}\cdot\nabla f=-(f-f_0)/\tau\). Approximation: relaxation-time approximation. Figure to show: streaming plus scattering toward equilibrium.
\end{notebox}

\begin{notebox}
\textbf{Slide 5: Fourier as a limiting model.}
Physical question: How does Fourier conduction emerge? Governing idea: when \(\mfp\ll L\), many scattering events erase directional memory, and the heat flux becomes local. Approximation: first-order near-equilibrium response to \(\nabla T\). Figure to show: many scattering events over one device length.
\end{notebox}

\begin{notebox}
\textbf{Slide 6: Why Cattaneo is not enough.}
Physical question: Does finite flux relaxation solve the nanoscale problem? Governing equation: \(\tau\partial\qvec/\partial t+\qvec=-k\nabla T\). Key limitation: it gives finite propagation speed but still lacks boundary-originating ballistic carriers. Figure to show: thermal wave picture versus particle flight picture.
\end{notebox}

\begin{notebox}
\textbf{Slide 7: Ballistic-diffusive split.}
Physical question: How can we keep ballistic memory without solving full phase space? Governing split: \(I_\omega=I_{b\omega}+I_{m\omega}\). Approximation: treat boundary-originating carriers by characteristics and scattered carriers by diffusion/P1 closure. Figure to show: carriers from wall versus carriers scattered inside the domain.
\end{notebox}

\begin{notebox}
\textbf{Slide 8: Final ballistic-diffusive PDE.}
Physical question: What equation will the FEM solve? Governing equation: medium energy obeys a Cattaneo-like PDE forced by \(-\nabla\cdot\qvec_b\). Key message: the ballistic term is the physical memory missing from Fourier and Cattaneo models. Figure to show: PDE box with \(u_m\), plus ray-based computation of \(\qvec_b\).
\end{notebox}

\begin{notebox}
\textbf{Slide 9: Thin-film benchmark.}
Physical question: How do we validate the implementation? Governing parameter: \(\mathrm{Kn}=\mfp/L\). Approximation being tested: reduced ballistic-diffusive equation versus Fourier and Cattaneo limits. Figure to show: profiles for \(\mathrm{Kn}=0.01,0.1,1,10\).
\end{notebox}

\begin{notebox}
\textbf{Slide 10: Relevance to superconducting qubits.}
Physical question: Why does this matter for qubit poisoning? Governing idea: energy deposited in a small region can move through phonons and other carriers before local equilibrium is established. Design relevance: heat sinks and traps should be evaluated by how they alter energy and quasiparticle transport near the junction-sensitive region.
\end{notebox}

\section{Module 1 implementation checklist}
Before moving to Module 2, the project should satisfy the following Module 1 checks:
\begin{enumerate}[leftmargin=1.6em]
    \item The code can compute \(q_b^*(\eta,t^*)\) with the correct finite-arrival lower limit \(\mu_t=\eta/(\mathrm{Kn}t^*)\).
    \item The code can solve Eq.~\eqref{eq:bd_1d_dimensionless_expanded} for several \(\mathrm{Kn}\) values.
    \item The diffusive limit trends toward Fourier-like behavior for \(\mathrm{Kn}\ll 1\) after early transients.
    \item The quasi-ballistic cases show delayed energy arrival and boundary-memory effects.
    \item The Cattaneo comparison is included to show why finite propagation speed alone is not the full physics.
    \item The plotted quantity is clearly labeled as equivalent internal-energy temperature, not blindly treated as local thermodynamic temperature.
\end{enumerate}

\section{References}
\begin{sloppypar}
\begin{enumerate}[label={[\arabic*]}]
\item G. Chen, ``Ballistic-Diffusive Heat-Conduction Equations,'' \emph{Physical Review Letters} \textbf{86}, 2297--2300 (2001).
\item G. Chen, ``Ballistic-Diffusive Equations for Transient Heat Conduction From Nano to Macroscales,'' \emph{Journal of Heat Transfer} \textbf{124}, 320--328 (2002).
\item D. D. Joseph and L. Preziosi, ``Heat waves,'' \emph{Reviews of Modern Physics} \textbf{62}, 375--391 (1990).
\item A. A. Joshi and A. Majumdar, ``Transient ballistic and diffusive phonon heat transport in thin films,'' \emph{Journal of Applied Physics} \textbf{74}, 31--39 (1993).
\item J. M. Ziman, \emph{Electrons and Phonons}, Oxford University Press (1960).
\item M. F. Modest, \emph{Radiative Heat Transfer}, McGraw-Hill (1993).
\end{enumerate}
\end{sloppypar}

\end{document}
